% ============================================================
%  TRABAJO ESPECIAL DE GRADO
%  Universidad Nororiental Privada Gran Mariscal de Ayacucho
%  Escuela de Ingeniería en Informática – Núcleo Bolívar
%  Normas: Manual de Elaboración de TEG – Ing. Patricia Torres
%
%  TÍTULO:
%  Sistema de digitalización y registros de productos con uso
%  de la Inteligencia Artificial para control de inventarios
%  en el Bodegón Pa' Donde Natty, Ciudad Bolívar, Estado Bolívar
%
%  ÁREA DE CONOCIMIENTO: Ingeniería de Software
%  LÍNEA DE INVESTIGACIÓN: 1.1 Sistemas de Información de apoyo
%  a la gestión de empresas comerciales y de servicios, de
%  pequeñas, mediana y gran envergadura.
% ============================================================

\documentclass[12pt, letterpaper]{report}

% --- Paquetes ---
\usepackage[spanish]{babel}
\usepackage[utf8]{inputenc}
\usepackage[T1]{fontenc}
\usepackage{times}
\usepackage{setspace}
\usepackage{geometry}
\usepackage{indentfirst}
\usepackage{titlesec}
\usepackage{fancyhdr}
\usepackage{hyperref}
\usepackage{textcase}
\usepackage{booktabs}
\usepackage{array}
\usepackage{tabularx}
\usepackage{multirow}
\usepackage{longtable}
\usepackage{graphicx}
\usepackage{float}
\usepackage{tocloft}
\usepackage{enumitem}

% --- Márgenes ---
\geometry{left=4cm, right=3cm, top=3cm, bottom=3cm}

% --- Configuración de índices al estilo del Manual ---
\renewcommand{\contentsname}{ÍNDICE GENERAL}
\renewcommand{\listfigurename}{ÍNDICE DE FIGURAS}
\renewcommand{\listtablename}{ÍNDICE DE TABLAS}

% --- Hyperref sin colores (impresión académica) ---
\hypersetup{hidelinks}

% --- Interlineado 1.5 ---
\setstretch{1.5}

% --- Sangría ---
\setlength{\parindent}{1.25cm}
\setlength{\parskip}{0pt}
\emergencystretch=2em

% --- Numeración romana de capítulos (exigida por el Manual) ---
\renewcommand{\thechapter}{\Roman{chapter}}
\renewcommand{\thesection}{\arabic{section}}
\renewcommand{\thesubsection}{\arabic{section}.\arabic{subsection}}

% --- Formato capítulo ---
\titleformat{\chapter}[display]
  {\normalfont\fontsize{14}{16}\bfseries\centering}
  {\MakeTextUppercase{\chaptertitlename}\ \thechapter}
  {6pt}
  {\fontsize{14}{16}\bfseries\centering\MakeTextUppercase}
\titlespacing*{\chapter}{0pt}{2cm}{24pt}

% --- Subtítulos ---
\titleformat{\section}
  {\normalfont\fontsize{13}{15}\bfseries}
  {\thesection}{1em}{}
\titlespacing*{\section}{0pt}{18pt}{6pt}

\titleformat{\subsection}
  {\normalfont\fontsize{12}{14}\bfseries}
  {\thesubsection}{1em}{}
\titlespacing*{\subsection}{0pt}{12pt}{4pt}


% --- Comando auxiliar para marcadores de figura ---
\newcommand{\figuraplaceholder}[1]{%
  \begin{center}
  \fbox{\parbox[c][4.5cm][c]{0.8\textwidth}{\centering
  \textit{[Espacio reservado para la figura]}\\[4pt]#1}}
  \end{center}}

% --- Comando para insertar figuras (imagen + título + fuente) ---
% Uso: \figuratesis{archivo}{ancho}{Figura N. Título descriptivo.}
\newcommand{\figuratesis}[3]{%
  \par\vspace{8pt}
  \phantomsection
  \addcontentsline{lof}{figure}{#3}%
  \noindent\begin{minipage}{\linewidth}
  \centering
  \includegraphics[width=#2\textwidth]{#1}\\[4pt]
  {\setstretch{1.0}\small\textbf{#3}\par}
  \end{minipage}
  \par\vspace{8pt}}

% --- Numeración de página ---
\pagestyle{fancy}
\fancyhf{}
\fancypagestyle{plain}{\fancyhf{}\rhead{\thepage}\renewcommand{\headrulewidth}{0pt}}
\setlength{\headheight}{14.5pt}
\rhead{\thepage}
\renewcommand{\headrulewidth}{0pt}

% ============================================================
\begin{document}
\raggedbottom

% ============================================================
%  PÁGINAS PRELIMINARES
% ============================================================

% --- Numeración en romanos minúsculos para preliminares ---
\pagenumbering{roman}
\hypersetup{pageanchor=false}
\pagestyle{empty}

% -------------------------------------------------------
%  PORTADA  (3 cm por los cuatro bordes, según el Manual)
% -------------------------------------------------------
\newgeometry{left=3cm, right=3cm, top=3cm, bottom=3cm}
\begin{titlepage}
\begin{center}

{\fontsize{12}{14}\selectfont\bfseries
UNIVERSIDAD NORORIENTAL PRIVADA\\[2pt]
``GRAN MARISCAL DE AYACUCHO''\\[2pt]
FACULTAD DE INGENIERÍA\\[2pt]
ESCUELA DE INGENIERÍA EN INFORMÁTICA\\[2pt]
NÚCLEO BOLÍVAR -- ESTADO BOLÍVAR\par}

\vspace{1.2cm}

% Logo institucional
\includegraphics[width=3cm]{uploads/logo_ugma.png}

\vspace{2.5cm}

{\fontsize{14}{18}\selectfont\bfseries
SISTEMA DE DIGITALIZACIÓN Y REGISTROS DE PRODUCTOS CON USO DE
LA INTELIGENCIA ARTIFICIAL PARA CONTROL DE INVENTARIOS EN EL
BODEGÓN PA' DONDE NATTY, CIUDAD BOLÍVAR, ESTADO BOLÍVAR.\par}

\vspace{2cm}

{\setstretch{1.0}\fontsize{12}{14}\selectfont
Trabajo Especial de Grado presentado como requisito parcial
para optar al título de Ingeniero en Informática.\par}

\vspace{2.5cm}

\begin{minipage}{0.45\textwidth}
\setstretch{1.0}
\centering
\textbf{Realizado por:}\\
{\footnotesize Gabriel Riera}\\
C.I.: V--31.153.478
\end{minipage}\hfill
\begin{minipage}{0.45\textwidth}
\setstretch{1.0}
\centering
\textbf{Tutor Académico:}\\
Ing. Carlos Moreno\\[10pt]
\end{minipage}

\vfill

{\setstretch{1.0}\fontsize{12}{14}\selectfont
Ciudad Bolívar, Mayo 2026.\par}

\end{center}
\end{titlepage}

% Restaurar márgenes del Manual para el resto del documento
\restoregeometry

% -------------------------------------------------------
%  PÁGINA DEL RESUMEN  (Manual: margen superior 3 cm)
% -------------------------------------------------------
\clearpage
\begin{center}
{\setstretch{1.0}\fontsize{12}{14}\selectfont\bfseries
UNIVERSIDAD NORORIENTAL PRIVADA\\
``GRAN MARISCAL DE AYACUCHO''\\
FACULTAD DE INGENIERÍA\\
ESCUELA DE INGENIERÍA EN INFORMÁTICA\\
NÚCLEO BOLÍVAR -- ESTADO BOLÍVAR\par}
\end{center}

\vspace{3pt}

{\setstretch{1.0}\noindent
\textbf{Línea de Investigación:} Sistemas de Información de apoyo a
la gestión de empresas comerciales y de servicios, de pequeñas,
mediana y gran envergadura.\\
\textbf{Área de Conocimiento:} Ingeniería de Software.\par}

\vspace{5pt}

\begin{center}
{\fontsize{12}{14}\selectfont\bfseries TÍTULO\par}
\vspace{4pt}
{\setstretch{1.0}\fontsize{12}{14}\selectfont\bfseries
SISTEMA DE DIGITALIZACIÓN Y REGISTROS DE PRODUCTOS CON USO DE LA
INTELIGENCIA ARTIFICIAL PARA CONTROL DE INVENTARIOS EN EL
BODEGÓN PA' DONDE NATTY, CIUDAD BOLÍVAR, ESTADO BOLÍVAR.\par}
\end{center}

\vspace{4pt}

{\setstretch{1.0}\noindent
\textbf{Autor(es):} Gabriel Riera\\
\textbf{C.I.:} V--31.153.478\\
\textbf{Tutor(es):} Ing.\ Carlos Moreno\\
\textbf{Lapso:} 1--2026\par}

\vspace{5pt}

\begin{center}
{\fontsize{12}{14}\selectfont\bfseries
RESUMEN DEL TRABAJO ESPECIAL DE GRADO\par}
\end{center}

\vspace{2pt}

{\setstretch{0.98}\noindent
El presente Trabajo Especial de Grado tuvo como propósito desarrollar
un prototipo funcional de un sistema de digitalización y registro de
productos con uso de Inteligencia Artificial para el control de
inventarios en el Bodegón Pa' Donde Natty C.A., ubicado en Ciudad
Bolívar, estado Bolívar. La investigación fue de nivel descriptivo,
con diseño de campo y apoyo documental; la población estuvo conformada
por los cinco (05) trabajadores de la empresa, abordados mediante censo
poblacional. Se aplicaron observación directa, entrevista no
estructurada, cuestionario, análisis documental y pruebas de caja
negra. El diagnóstico evidenció discrepancias entre el inventario
físico y los registros, además de ciclos de conteo de 4 a 6 horas. El
prototipo, desarrollado con Python/FastAPI, Next.js y
PostgreSQL/Supabase, integró un modelo multimodal mediante API y
alcanzó 93\,\% de exactitud de reconocimiento, 2,4 s de tiempo de
respuesta promedio y 100\,\% de cumplimiento funcional. Se concluye
que la solución propuesta es técnica, operativa y económicamente
viable para digitalizar el inventario en empresas comerciales de
pequeña y mediana envergadura.\par}

\vspace{3pt}

{\setstretch{0.98}\noindent
\textbf{Descriptores:} Inteligencia Artificial, control de
inventarios, reconocimiento de imágenes, digitalización, sistemas de
información.\par}

\vspace{0.15cm}

\noindent\begin{minipage}{0.45\textwidth}
\centering
Tutor Académico\\[8pt]
\rule{5cm}{0.4pt}\\
(Firma)
\end{minipage}\hfill
\begin{minipage}{0.45\textwidth}
\centering
Estudiante\\[8pt]
\rule{5cm}{0.4pt}\\
(Firma)
\end{minipage}

% -------------------------------------------------------
%  DEDICATORIA
% -------------------------------------------------------
\clearpage
\begin{flushright}
{\fontsize{14}{16}\selectfont\bfseries DEDICATORIA\par}
\end{flushright}
\vspace{0.6cm}

\begin{flushright}
\begin{minipage}{0.72\textwidth}
\footnotesize
\setstretch{1.12}
\itshape
A Dios todopoderoso, por darme la vida, la fortaleza y la sabiduría necesarias para alcanzar esta meta, y por sostenerme en los momentos en que sentí que no podía continuar.

A mi padre, Gabriel, quien físicamente ya no está conmigo pero vive cada día en mi memoria y en cada uno de mis logros. Este triunfo es tuyo tanto como mío, papá. Sé que desde donde estás me acompañas y celebras conmigo. Todo lo que soy y todo lo que haga, siempre será en tu honor.

A mi madre, Maylin, mi pilar, mi refugio y mi más grande ejemplo de fortaleza. Por sacar adelante a nuestra familia con tanto amor y tanto sacrificio, por no soltarme nunca la mano y por enseñarme que con esfuerzo y fe todo es posible. Este logro es la cosecha de cada una de tus desveladas. Te amo, mamá.

A mis abuelos, Guadalupe y Marcos, mis segundos padres. Por la crianza, por el amor, por las enseñanzas y por estar presentes en cada etapa de mi vida cuando más los necesité. No tengo palabras para agradecer todo lo que han hecho por mí; este logro también lleva sus nombres.

A mis hermanas, Camilla y Tatiana, mis compañeras de vida. Que este logro les sirva como ejemplo de que con disciplina y constancia los sueños se alcanzan. Las amo.

A mi novia, Jhosmeli, por ser mi apoyo, mi calma y mi compañía en este camino. Gracias por entender cada hora de estudio, cada ausencia y por creer en mí incluso cuando yo mismo dudé. Este logro también es tuyo.

A mi tía Gabriela (Gaby) y a mi prima Valeria, por su cariño, su compañía y por estar siempre presentes a lo largo de este camino.

A toda mi familia, por su amor incondicional y por ser el cimiento sobre el que construyo cada paso.
\end{minipage}
\end{flushright}

% -------------------------------------------------------
%  AGRADECIMIENTO
% -------------------------------------------------------
\clearpage
\begin{center}
{\fontsize{14}{16}\selectfont\bfseries AGRADECIMIENTO\par}
\end{center}
\vspace{0.8cm}

{\small
\setstretch{1.15}
\itshape
En primer lugar, a Dios todopoderoso, por la salud, la sabiduría y la perseverancia que me regaló durante todo este proceso. Por poner en mi camino a las personas correctas en los momentos exactos, y por ser la luz que iluminó cada etapa de esta investigación.

A mi padre Gabriel, que aunque ya no está físicamente a mi lado, me dejó la mejor herencia que un padre puede dejar: valores, principios y el ejemplo de un hombre trabajador. Gracias, papá, porque sé que desde el cielo me cuidaste durante todo este proceso.

A mi madre, Maylin, mi infinita gratitud por su amor incondicional, su esfuerzo diario y por nunca dejar de creer en mí. Cada palabra de aliento suya, cada sacrificio silencioso, cada oración, hicieron posible que hoy esté culminando esta etapa. Gracias, mamá, por ser madre y padre al mismo tiempo.

A mis abuelos, Guadalupe y Marcos, por su amor sin medida, por sus consejos, por su compañía y por ese apoyo que nunca me faltó. Gracias por criarme, por enseñarme y por estar siempre que los necesité; este logro es también el fruto de su crianza.

A mis hermanas, Camilla y Tatiana, por estar siempre presentes, por las risas, los buenos momentos y por el cariño que nos une. Gracias por ser parte fundamental de mi vida.

A mi novia, Jhosmeli, gracias por la paciencia, por escucharme hablar incontables veces de este trabajo, por acompañarme en las noches de desvelo, por animarme cuando el camino se hacía cuesta arriba y por celebrar cada pequeño avance como si fuera tuyo. Tu compañía hizo este proceso mucho más llevadero.

A mi tía Gabriela (Gaby), por su cariño constante, su apoyo y por estar siempre pendiente de mis estudios; y a mi prima Valeria, por su compañía y su afecto a lo largo de estos años.

A toda mi familia, tíos, primos y demás seres queridos, por su apoyo y por cada palabra de aliento que recibí durante estos años de estudio.

A la Universidad Nororiental Privada Gran Mariscal de Ayacucho, mi casa de estudios, por abrirme sus puertas y permitirme formarme como profesional en la Escuela de Ingeniería en Informática del Núcleo Bolívar. A cada uno de los profesores que a lo largo de la carrera compartieron conmigo sus conocimientos y experiencias, mi más sincero reconocimiento.

De manera muy especial, agradezco a mi tutor académico, el Ing. Carlos Moreno, por su orientación constante, su disposición permanente, su paciencia y por cada recomendación oportuna que dio forma y rigor a esta investigación. Su acompañamiento fue determinante para alcanzar los resultados aquí presentados.

A la coordinación de TEG-Pasantía, en especial a la Ing. Patricia Torres, por su guía y por velar por el cumplimiento de las normas que orientan la elaboración del Trabajo Especial de Grado.

Al Bodegón Pa' Donde Natty C.A., en la persona de su gerencia y de todo su personal, por abrirme las puertas de la empresa, brindarme la información necesaria y prestar su tiempo y disposición durante la fase de diagnóstico y validación del prototipo. Sin su colaboración, este trabajo no habría sido posible.

A mis compañeros de carrera y amigos, con quienes compartí jornadas de estudio, trabajos en equipo y muchas horas frente al computador. Gracias por hacer del camino algo más llevadero.

A todas aquellas personas que, directa o indirectamente, contribuyeron a la culminación de este Trabajo Especial de Grado, mi más sincero y profundo agradecimiento.
\par}

% -------------------------------------------------------
%  ÍNDICE GENERAL
% -------------------------------------------------------
\clearpage
\begin{center}
{\fontsize{14}{16}\selectfont\bfseries ÍNDICE GENERAL\par}
\end{center}
\vspace{8pt}

\begingroup
\setstretch{1.0}
\renewcommand{\contentsname}{}
\vspace{-2.5cm}
\tableofcontents
\endgroup

% -------------------------------------------------------
%  ÍNDICE DE TABLAS
% -------------------------------------------------------
\clearpage
\begin{center}
{\fontsize{14}{16}\selectfont\bfseries ÍNDICE DE TABLAS\par}
\end{center}
\vspace{8pt}

\begingroup
\setstretch{1.0}
\renewcommand{\listtablename}{}
\vspace{-2.5cm}
\listoftables
\endgroup

% -------------------------------------------------------
%  ÍNDICE DE FIGURAS
% -------------------------------------------------------
\clearpage
\begin{center}
{\fontsize{14}{16}\selectfont\bfseries ÍNDICE DE FIGURAS\par}
\end{center}
\vspace{8pt}

\begingroup
\setstretch{1.0}
\renewcommand{\listfigurename}{}
\vspace{-2.5cm}
\listoffigures
\endgroup

% -------------------------------------------------------
%  ÍNDICE DE ANEXOS
% -------------------------------------------------------
\clearpage
\begin{center}
{\fontsize{14}{16}\selectfont\bfseries ÍNDICE DE ANEXOS\par}
\end{center}
\vspace{12pt}

\begingroup
\setstretch{1.0}
\noindent
\begin{tabular}{@{}p{2.2cm}p{10cm}r@{}}
\textbf{Anexo A.} & Cuestionario aplicado al personal del Bodegón Pa' Donde Natty C.A. & 1 \\[6pt]
\textbf{Anexo B.} & Guion de entrevista no estructurada dirigida al gerente y al encargado de almacén. & 2 \\[6pt]
\textbf{Anexo C.} & Guía de observación estructurada del proceso de inventario. & 3 \\[6pt]
\textbf{Anexo D.} & Lista de cotejo para la validación funcional del prototipo. & 4 \\[6pt]
\textbf{Anexo E.} & Especificación de Requerimientos de Software (SRS) consolidada. & 5 \\[6pt]
\end{tabular}
\endgroup

% -------------------------------------------------------
%  INTRODUCCIÓN
% -------------------------------------------------------
\clearpage
\pagenumbering{arabic}
\hypersetup{pageanchor=true}
\setcounter{page}{1}
\pagestyle{fancy}

\begin{center}
{\fontsize{14}{16}\selectfont\bfseries INTRODUCCIÓN\par}
\end{center}
\vspace{12pt}
\addcontentsline{toc}{chapter}{INTRODUCCIÓN}

\hspace{1.25cm}En la actualidad, la transformación digital de los
procesos empresariales se ha consolidado como uno de los pilares de la
competitividad organizacional. La adopción de tecnologías avanzadas
---en particular la Inteligencia Artificial--- ha modificado la forma
en que las empresas, sin importar su tamaño, gestionan sus
operaciones, automatizan tareas y soportan la toma de decisiones. Sin
embargo, este avance no ha permeado con la misma intensidad a todos
los sectores: en el comercio de consumo masivo de pequeña y mediana
escala, especialmente en economías con limitaciones tecnológicas, una
parte importante de las empresas continúa operando con métodos
manuales o semi-digitales que generan pérdidas recurrentes por
errores de registro, desabastecimiento y vencimiento de mercancía.

El presente Trabajo Especial de Grado se inscribe en este contexto.
Tiene como propósito desarrollar un prototipo funcional de un sistema
de digitalización y registros de productos con uso de Inteligencia
Artificial para el control de inventarios en el Bodegón Pa' Donde
Natty C.A., empresa comercial ubicada en Ciudad Bolívar, estado
Bolívar. La pertinencia del estudio radica en que combina, por un
lado, una necesidad concreta y verificable de una empresa real ---la
sustitución de un proceso manual propenso a errores--- y, por otro,
una oportunidad tecnológica de actualidad: el uso de modelos de
lenguaje multimodales consumidos mediante su interfaz de programación
de aplicaciones (API), que permite reconocer productos a partir de su
imagen sin necesidad de entrenar un modelo propio.

Desde el punto de vista conceptual, la investigación se sustenta en
nociones del control de inventarios, la digitalización de procesos
empresariales, la Inteligencia Artificial, el reconocimiento de
imágenes, la ingeniería de software y las bases de datos relacionales,
articuladas a partir de los aportes de autores como Ballou (2004),
Guerrero (2009), Laudon y Laudon (2016), Russell y Norvig (2016),
Sommerville (2011) y Date (2001), entre otros. Desde el punto de vista
metodológico, el estudio se aborda como una investigación de tipo
descriptivo, con diseño de campo apoyado en fuentes documentales, en
la que se aplican técnicas como la observación directa, la entrevista
no estructurada, el cuestionario, el análisis documental y la
observación sistemática mediante pruebas de software.

El informe se estructura en cinco (05) capítulos. El Capítulo I,
\textit{El Problema}, describe el planteamiento, los objetivos
---general y específicos---, la justificación, el alcance, las
limitaciones y la línea de investigación en que se enmarca el trabajo.
El Capítulo II, \textit{Marco Teórico}, presenta los antecedentes, los
fundamentos teóricos, las bases legales y la referencia conceptual que
sirven de soporte académico al estudio. El Capítulo III, \textit{Marco
Metodológico}, expone el tipo y diseño de la investigación, la
población y muestra, las técnicas e instrumentos de recolección de
datos, el cuadro de operacionalización de la variable, el
procedimiento por objetivo y las técnicas de análisis. El Capítulo IV,
\textit{Análisis de Resultados}, expone, organiza e interpreta los
hallazgos correspondientes al diagnóstico, al diseño, al desarrollo y
a la validación del prototipo, e incluye un breve análisis de
factibilidad. Finalmente, el Capítulo V, \textit{Conclusiones y
Recomendaciones}, sintetiza los logros del trabajo en relación con
cada objetivo específico y plantea sugerencias para la implementación,
el mantenimiento y la extensión del sistema. Al cierre del documento
se incluyen las referencias bibliográficas consultadas y los anexos
que respaldan el proceso investigativo.

% ============================================================


% -------------------------------------------------------
%  CAPÍTULO I  –  EL PROBLEMA
% -------------------------------------------------------
\chapter{El Problema}

En el presente capítulo se describe de manera amplia la situación
objeto de estudio, ubicándola en un contexto que permite comprender
su origen, sus relaciones y las incógnitas por responder. Se parte de
una visión general hacia lo particular y se desarrollan el
planteamiento del problema, los objetivos de la investigación, la
justificación, el alcance y las limitaciones del estudio.

% -------------------------------------------------------
\section{Planteamiento del Problema}
% -------------------------------------------------------

En la actualidad, la gestión eficiente del inventario representa uno de
los pilares fundamentales para la competitividad y sostenibilidad de las
empresas comerciales. A escala global, la transformación digital ha
impulsado la adopción de tecnologías avanzadas, como la Inteligencia
Artificial (IA) y los sistemas de digitalización de registros, con el
propósito de optimizar los procesos logísticos, reducir costos
operativos y mejorar la toma de decisiones en tiempo real. Organizaciones
de todos los sectores han migrado de los métodos tradicionales basados en
papel hacia plataformas automatizadas que ofrecen mayor precisión,
trazabilidad y velocidad en el manejo de la información.

En América Latina, esta tendencia global ha comenzado a permear
gradualmente el tejido empresarial, aunque con brechas importantes según
el nivel de desarrollo tecnológico de cada país. Naciones como Brasil,
México, Colombia y Chile han avanzado de manera significativa en la
adopción de sistemas digitales de gestión de inventarios, impulsados por
un ecosistema tecnológico en expansión y el crecimiento del comercio
electrónico regional. Sin embargo, en buena parte de la región,
especialmente en economías con mayor volatilidad y menor acceso a
financiamiento tecnológico, las micro, pequeñas y medianas empresas
(mipymes) del sector comercial ---particularmente las empresas
distribuidoras de consumo masivo y bodegas de alimentos--- continúan
operando con métodos manuales o semi-digitales que las exponen a errores
como pérdidas por vencimiento de productos, desabastecimiento de
artículos de alta rotación, duplicidad de registros al transcribir datos
y discrepancias entre el inventario físico y el digital. Organismos como
la Comisión Económica para América Latina y el Caribe (CEPAL) han
señalado que la transformación digital de las mipymes constituye uno de
los desafíos estructurales más urgentes de la región para elevar su
productividad y reducir la brecha tecnológica frente a empresas de mayor
escala.

En Venezuela, y en particular en la región de Guayana, pese a las
condiciones económicas adversas, el sector comercial ha buscado
paulatinamente modernizar sus operaciones. Sin embargo, gran parte de
las pequeñas y medianas empresas (pymes), en especial las bodegas y
distribuidoras de productos de consumo masivo ubicadas en zonas de mayor
vulnerabilidad económica como el estado Bolívar, aún dependen de
registros manuales o sistemas desactualizados. Esta situación genera
consecuencias directas sobre su operatividad: pérdidas de productos por
vencimiento no detectado a tiempo, desabastecimiento recurrente de
artículos de alta rotación, registros duplicados o incorrectos al
transcribir datos manualmente, y una notoria discrepancia entre el
inventario físico real y el que reflejan los libros de control. Tales
problemas limitan la competitividad de estas empresas y las exponen a
pérdidas económicas sistemáticas que podrían evitarse con una adecuada
gestión tecnológica.

En el municipio Heres del estado Bolívar, específicamente en Ciudad
Bolívar, parroquia Agua Salada, urbanización Los Próceres, opera el
Bodegón Pa' Donde Natty C.A., empresa del sector
comercial dedicada a la distribución y venta de productos de consumo
masivo, entre los que se encuentran alimentos no perecederos (arroz,
harina, azúcar, aceite, caraotas, entre otros), bebidas (agua mineral,
refrescos, jugos), artículos de higiene personal y productos de limpieza
del hogar. El control de existencias de esta amplia gama de productos se
lleva actualmente mediante procedimientos manuales: conteos físicos
periódicos, registros en hojas de cálculo y apuntes en papel. Este
esquema operativo ha derivado en errores concretos y recurrentes, tales
como: discrepancias entre el inventario físico y el digital,
desabastecimiento de productos de alta rotación sin detección oportuna,
pérdidas por vencimiento de mercancía no identificada a tiempo, y
retardos en la reposición de stock. Todo ello impide una gestión
dinámica ---entendida como la capacidad de actualizar el registro de
existencias en tiempo real ante cada movimiento de mercancía--- y
confiable ---es decir, que garantice la exactitud e integridad de los
datos registrados en todo momento---, elementos indispensables para la
toma de decisiones en un negocio de alta rotación como el que nos
ocupa.

De la situación descrita surgen las siguientes interrogantes que orientan
la presente investigación:

\begin{enumerate}
  \item ¿Cuáles son las necesidades funcionales y técnicas del proceso
        actual de registro y control de inventarios del Bodegón Pa'
        Donde Natty C.A.?

  \item ¿Qué arquitectura lógica, modelo de datos e interfaz de usuario
        debe poseer el sistema de digitalización y registros de
        productos con IA para satisfacer los requerimientos
        identificados?

  \item ¿Cómo desarrollar el prototipo funcional integrando el
        reconocimiento de productos mediante un modelo de lenguaje
        multimodal de Inteligencia Artificial consumido a través de
        su API?

  \item ¿En qué medida el prototipo desarrollado cumple con los
        requerimientos definidos al ser sometido a pruebas en un
        entorno controlado?
\end{enumerate}

% -------------------------------------------------------
\section{Objetivos de la Investigación}
% -------------------------------------------------------

\subsection{Objetivo General}

Desarrollar un prototipo funcional de un sistema de digitalización y
registros de productos con uso de la Inteligencia Artificial para el
control de inventarios en el Bodegón Pa' Donde Natty C.A., Ciudad
Bolívar, estado Bolívar.

\subsection{Objetivos Específicos}

\begin{enumerate}
  \item Diagnosticar las necesidades funcionales y técnicas del proceso
        actual de registro y control de inventarios del Bodegón Pa'
        Donde Natty C.A., mediante la aplicación de técnicas de
        recolección de información.

  \item Diseñar la arquitectura lógica, el modelo de datos y la
        interfaz de usuario del sistema de digitalización y registros
        de productos con IA, consolidados en una Especificación de
        Requerimientos de Software (SRS, por sus siglas en inglés).

  \item Desarrollar el prototipo funcional del sistema integrando el
        reconocimiento de productos mediante un modelo de lenguaje
        multimodal de Inteligencia Artificial, consumido a través de su
        interfaz de programación de aplicaciones (API).

  \item Validar el funcionamiento del prototipo mediante pruebas en un
        entorno controlado, verificando el cumplimiento de los
        requerimientos definidos.
\end{enumerate}

% -------------------------------------------------------
\section{Justificación}
% -------------------------------------------------------

La presente investigación se justifica desde múltiples perspectivas que
resaltan su pertinencia académica, tecnológica y empresarial.

Desde el punto de vista \textbf{tecnológico}, la integración de la
Inteligencia Artificial en los procesos de inventario representa un
avance significativo en la modernización de las operaciones comerciales.
El uso de técnicas como el reconocimiento de imágenes, el aprendizaje
automático y el procesamiento de datos en tiempo real permite automatizar
tareas repetitivas, reducir el margen de error humano y generar
información estratégica para la toma de decisiones.

Desde la perspectiva \textbf{empresarial}, el Bodegón Pa' Donde Natty
C.A. se beneficiará directamente con la reducción de las pérdidas
generadas por desabastecimiento o exceso de inventario, problemática que,
de acuerdo con Johan et al.\ (2024), puede ser prevenida mediante la
aplicación de algoritmos de Inteligencia Artificial que predicen la
demanda y ajustan los niveles de stock de forma proactiva, evitando
tanto la rotura como la sobreacumulación de mercancía. La digitalización
de los registros también aportará transparencia y trazabilidad a todos
los movimientos del almacén, y abrirá la posibilidad de escalar las
operaciones de la empresa de forma sostenida.

En el ámbito \textbf{académico}, esta investigación aporta un modelo
replicable para otras pequeñas y medianas empresas de la región que
enfrentan problemáticas similares, contribuyendo al cuerpo de
conocimiento sobre la aplicación de la IA en contextos empresariales
venezolanos. Asimismo, representa un aporte concreto al área de sistemas
de información e Ingeniería de la Universidad Nororiental Privada Gran
Mariscal de Ayacucho.

Desde el plano \textbf{social}, el proyecto favorece la generación de
empleo calificado en el área tecnológica dentro del estado Bolívar y
fomenta la cultura digital en el sector comercial local, contribuyendo
al desarrollo económico de la región.

% -------------------------------------------------------
\section{Alcance}
% -------------------------------------------------------

La investigación comprende el análisis, diseño, desarrollo y validación
de un prototipo funcional de un sistema de digitalización y registros de
productos con uso de Inteligencia Artificial, orientado al control de
inventarios del Bodegón Pa' Donde Natty C.A., ubicado en Ciudad
Bolívar, estado Bolívar. El prototipo abarcará los módulos de: registro
digital de productos, reconocimiento automatizado mediante técnicas de
IA, gestión de stock (entradas, salidas y alertas de reorden) y
generación de reportes gerenciales. Su validación se realizará mediante
pruebas funcionales en un entorno controlado de desarrollo, sin que ello
implique la implantación definitiva del sistema en la operación real de
la empresa. El alcance del proyecto está delimitado al área de almacén e
inventario de la empresa mencionada, sin incluir módulos de facturación,
contabilidad ni comercio electrónico, los cuales podrán abordarse en
investigaciones futuras.

% -------------------------------------------------------
\section{Limitaciones}
% -------------------------------------------------------

Durante el desarrollo de la investigación se prevén las siguientes
limitaciones:

\begin{itemize}
  \item \textbf{Acceso a la información}: La disponibilidad y
        disposición del personal de la empresa para suministrar
        información detallada sobre sus procesos internos podría
        condicionar la profundidad del diagnóstico.

  \item \textbf{Dependencia de servicios externos}: el reconocimiento
        de productos se apoya en un servicio de Inteligencia Artificial
        de terceros (API de un modelo multimodal) y el prototipo se
        despliega en plataformas de capa gratuita (Vercel, Render y
        Supabase). Los límites de uso de dichas capas gratuitas ---así
        como la suspensión por inactividad del backend y de la base de
        datos--- condicionan el volumen de operaciones y pueden
        introducir una demora en la primera petición tras un período de
        inactividad.

  \item \textbf{Conectividad}: dado que el reconocimiento se realiza
        mediante una API en la nube, el prototipo requiere conexión a
        Internet durante su operación; la calidad del enlace puede
        influir en el tiempo de respuesta del reconocimiento.

  \item \textbf{Tiempo}: El período académico disponible para el
        desarrollo del Trabajo Especial de Grado limita el alcance de
        las pruebas y la validación exhaustiva del sistema en un entorno
        de producción real.
\end{itemize}

% -------------------------------------------------------
\section{Línea de Investigación}
% -------------------------------------------------------

El presente Trabajo Especial de Grado se enmarca dentro del
\textbf{Área de Conocimiento de Ingeniería de Software}, correspondiente
a la \textbf{Línea de Investigación 1.1: Sistemas de Información de
apoyo a la gestión de empresas comerciales y de servicios, de pequeñas,
mediana y gran envergadura}, según la clasificación establecida por la
Escuela de Ingeniería en Informática de la Universidad Nororiental
Privada Gran Mariscal de Ayacucho, Núcleo Bolívar. Esta línea
comprende el análisis, desarrollo y evaluación de programas orientados
a apoyar la gestión de empresas comerciales y de servicios mediante el
uso de lenguajes de programación, elaboración de algoritmos y
construcción de bases de datos, lo cual se corresponde plenamente con
el objeto y los objetivos del presente estudio.

% ============================================================

% ============================================================

% -------------------------------------------------------
%  CAPÍTULO II  –  MARCO TEÓRICO
% -------------------------------------------------------
\chapter{Marco Teórico}

El marco teórico representa la especificación de las referencias
teóricas que se utilizan para formular y desarrollar los argumentos
que sustentan la investigación. En este capítulo se desarrollan los
antecedentes de la investigación, los fundamentos teóricos, las bases
legales y la referencia conceptual relacionados con el problema
planteado.

% -------------------------------------------------------
\section{Antecedentes de la Investigación}
% -------------------------------------------------------

Los antecedentes recogen trabajos de investigación previos que guardan
relación directa con el problema planteado y aportan bases metodológicas
y técnicas al presente estudio. Se presentan en orden cronológico del
más reciente al más antiguo, con un máximo de tres (03) según las
normas de la institución.

En el año 2024, Johan, S., Riascos-Guerrero, J. A., Galván-Colonia, E.
y Pincay-Lozada, J. L., en su artículo científico titulado:
\textbf{``Estrategias basadas en Inteligencia Artificial para la gestión
de inventarios en la cadena de suministro''}, publicado en la Revista
\textit{Tecnología en Marcha}, Vol. 37, N° 6, págs. 88--97 (DOI:
10.18845/tm.v37i6.7271), de la Universidad Cooperativa de Colombia,
presentado en el XI Congreso Internacional en Inteligencia Ambiental,
Ingeniería de Software, Salud Electrónica y Móvil (AmITIC). Definiendo
esta como una investigación de tipo documental--analítica, de enfoque
cuantitativo. El objetivo general fue analizar y sistematizar las
estrategias basadas en Inteligencia Artificial aplicadas a la gestión
de inventarios en la cadena de suministro empresarial. Entre sus
principales hallazgos, los autores concluyeron que la aplicación de
algoritmos avanzados de aprendizaje automático sobre datos históricos
de ventas, patrones de compra y condiciones económicas permite predecir
con alta precisión la demanda futura, lo que posibilita ajustar los
niveles de inventario de forma efectiva y evitar situaciones de escasez
o exceso de stock; asimismo, identificaron que la IA contribuye a
detectar cuellos de botella y a sugerir mejoras en logística y
distribución, elevando la eficiencia operativa y la satisfacción del
cliente. El mencionado trabajo aporta a la presente investigación el
marco conceptual sobre la aplicación estratégica de la IA en la gestión
de inventarios y los criterios para la selección de algoritmos de
predicción de demanda en contextos empresariales de pequeña y mediana
escala.

En el año 2023, Wong Portillo, L., De la Torre Ganoza, R. y Reyes
Herrera, A., en su trabajo de investigación titulado:
\textbf{``Aplicación de reconocimiento de imágenes para productos en el
sector retail''}, presentado en la Universidad Peruana de Ciencias
Aplicadas (UPC), Perú, disponible en el Repositorio Académico UPC
(handle 10757/667951). Definiendo esta como una investigación de tipo
aplicada y descriptiva. El objetivo general fue revisar y sistematizar
investigaciones sobre la detección de objetos mediante técnicas de Deep
Learning para, posteriormente, implementar un modelo propio de
reconocimiento usando YOLO con un dataset personalizado de productos
de retail. Los autores concluyeron que las técnicas de detección de
objetos basadas en aprendizaje profundo, especialmente los modelos
YOLO, constituyen la opción más adecuada para el reconocimiento
automatizado de productos en entornos comerciales, destacando su alta
velocidad de inferencia y precisión en tiempo real como ventajas
determinantes frente a métodos tradicionales de clasificación de
imágenes. El aporte de esta investigación al presente trabajo radica
en evidenciar la viabilidad y el valor del reconocimiento automatizado
de productos a partir de imágenes en entornos comerciales y de almacén;
si bien dicho antecedente recurre al entrenamiento de un modelo propio,
el presente trabajo adopta una alternativa más reciente ---el consumo
de un modelo de lenguaje multimodal mediante su API--- para alcanzar el
mismo fin de reconocer productos a partir de su imagen.

En el año 2023, Román Veliz, A. y Arce Ríos, M., en su Trabajo Especial
de Grado titulado: \textbf{``Implementación de un sistema de gestión de
inventarios para mejorar la eficiencia en la logística de
aprovisionamiento de la planta lechera `Concelac' en la ciudad de
Concepción-2022''}, presentado para optar el Título Profesional de
Ingeniero Industrial en la Universidad Continental, Huancayo, Perú,
disponible en el Repositorio Continental (handle 20.500.12394/13413).
Investigación de tipo aplicada, correlacional y de campo, con diseño
experimental-transversal. El objetivo general fue implementar un sistema
de gestión de inventarios mediante desarrollo de software para mejorar
la eficiencia en los procesos de compras, inventario y almacenaje de la
planta, cuya eficiencia de diagnóstico inicial era de 37,5\%. Los
autores concluyeron que la implementación del sistema elevó el nivel de
exactitud del inventario y redujo significativamente los costos de
almacenamiento, al tiempo que permitió la generación de alertas
automáticas de stock mínimo y dashboards de rotura de stock en tiempo
real, mejorando la toma de decisiones gerenciales. El aporte de este
antecedente a la presente investigación se centra en el modelo de diseño
de software para gestión de inventarios, los módulos de alertas de
reorden y el esquema de base de datos relacional empleado, los cuales
sirven de referencia directa para el desarrollo del sistema propuesto.

% -------------------------------------------------------
\section{Fundamentos Teóricos}
% -------------------------------------------------------

Los fundamentos teóricos comprenden el conjunto de conceptos,
proposiciones y teorías de diversos autores que constituyen el soporte
académico y técnico de la investigación, permitiendo ubicar el problema
en un enfoque teórico determinado y relacionar la teoría con el objeto
de estudio.

% -------
\subsection{Inventario y Control de Inventarios}
% -------

El inventario, según Ballou (2004), es definido como el conjunto de
artículos o recursos usados por una organización que pueden ser
almacenados para satisfacer una necesidad futura. En el contexto
empresarial, el inventario representa uno de los activos más
importantes y su gestión inadecuada puede derivar en pérdidas
significativas por obsolescencia, vencimiento o ruptura de stock.
Guerrero (2009) señala que el control de inventarios es el proceso
sistemático de supervisar, registrar y administrar las existencias de
productos, con el propósito de garantizar la disponibilidad de los
mismos en el momento y lugar requerido, minimizando los costos
asociados.

Los sistemas de control de inventarios se clasifican, según Chase,
Aquilano y Jacobs (2009), en dos grandes categorías: los sistemas de
revisión continua, donde el nivel de inventario se monitorea de forma
permanente, y los sistemas de revisión periódica, donde la
verificación se realiza en intervalos de tiempo fijos. La tendencia
actual apunta hacia la automatización de estos sistemas mediante el uso
de tecnologías digitales que permitan la revisión continua y en tiempo
real.

% -------
\subsection{Digitalización de Procesos Empresariales}
% -------

La digitalización, según Laudon y Laudon (2016), consiste en la
conversión de información y procesos analógicos a formatos digitales
procesables por sistemas computacionales. En el ámbito empresarial, la
digitalización de registros transforma los flujos de trabajo manuales
en procesos automatizados, reduciendo el error humano, aumentando la
velocidad de procesamiento y facilitando el acceso centralizado a la
información. Schwab (2016) destaca que la digitalización constituye
el primer escalón de la transformación digital, siendo condición
necesaria para la posterior incorporación de tecnologías avanzadas
como la Inteligencia Artificial.

Para empresas del sector comercial como el Bodegón Pa' Donde Natty
C.A., la digitalización de los registros de inventario implica la
sustitución de hojas de cálculo y apuntes manuales por un sistema de
base de datos centralizado, accesible desde múltiples dispositivos y
capaz de registrar cada movimiento de productos en tiempo real, con
trazabilidad completa de las transacciones.

% -------
\subsection{Inteligencia Artificial}
% -------

La Inteligencia Artificial (IA), según Russell y Norvig (2016), es la
rama de las ciencias de la computación dedicada al diseño de sistemas
capaces de realizar tareas que, cuando son ejecutadas por seres humanos,
requieren inteligencia. Estos sistemas aprenden de la experiencia, se
adaptan a nuevas situaciones y ejecutan tareas cognitivas como el
reconocimiento de patrones, la toma de decisiones y la predicción.

Dentro de la IA, el \textbf{aprendizaje automático} (machine learning)
ocupa un lugar central. Mitchell (1997) lo define como el campo de
estudio que confiere a los computadores la capacidad de aprender sin
ser explícitamente programados. En el contexto de la gestión de
inventarios, el aprendizaje automático permite construir modelos
predictivos de demanda, detectar anomalías en el stock y automatizar
las decisiones de reorden, reduciendo la dependencia de criterios
subjetivos del operario. Johan et al.\ (2024) confirman que la
aplicación de estos algoritmos sobre grandes volúmenes de datos
históricos de ventas logra una predicción de la demanda futura con
alta precisión, posibilitando el ajuste proactivo de los niveles de
inventario.

% -------
\subsection{Modelos de Lenguaje Multimodales}
% -------

Los modelos de lenguaje multimodales (MLLM, por sus siglas en inglés)
constituyen una evolución reciente de la Inteligencia Artificial y, en
particular, del aprendizaje automático descrito por Mitchell (1997) y
del enfoque general de los sistemas inteligentes planteado por Russell
y Norvig (2016). A diferencia de los modelos tradicionales, que
procesan un único tipo de dato, un modelo multimodal es capaz de
recibir y relacionar simultáneamente distintas modalidades de
información ---texto e imagen, principalmente--- y producir una
respuesta en lenguaje natural. Esta capacidad permite que, a partir de
la fotografía de un producto y de una instrucción textual, el modelo
identifique el artículo y describa sus atributos sin necesidad de un
entrenamiento específico por parte del desarrollador.

En el contexto del presente trabajo, el modelo de lenguaje multimodal
es el componente que realiza el reconocimiento de los productos del
Bodegón Pa' Donde Natty C.A. a partir de la imagen capturada,
sustituyendo el enfoque tradicional de entrenar un modelo propio de
detección de objetos por el consumo de un servicio de IA ya entrenado
y de propósito general.

% -------
\subsection{Reconocimiento de Imágenes mediante Inteligencia Artificial}
% -------

El reconocimiento de imágenes es la capacidad de un sistema
computacional para identificar y clasificar el contenido de una imagen
digital, asignando etiquetas o categorías a los elementos visuales
presentes en ella. Históricamente esta tarea requería construir y
entrenar modelos especializados sobre conjuntos de imágenes
etiquetadas; los trabajos previos en el área, como el de Wong Portillo
et al.\ (2023), demostraron la viabilidad y el valor de reconocer
automáticamente productos a partir de su imagen en entornos
comerciales. Con la aparición de los modelos de lenguaje multimodales,
el reconocimiento de imágenes puede realizarse mediante un servicio de
IA de propósito general, sin que el desarrollador deba construir un
conjunto de datos ni entrenar un modelo, lo que reduce
significativamente la complejidad técnica y el costo del proyecto.

Para el sistema propuesto, el reconocimiento de imágenes es la función
central del módulo de registro: a partir de la foto del producto se
obtiene su identificación, que alimenta la actualización automática del
inventario.

% -------
\subsection{Servicios de IA mediante API e Ingeniería de Prompts}
% -------

Una interfaz de programación de aplicaciones (API, por sus siglas en
inglés) es, de acuerdo con los principios de la ingeniería de software
descritos por Sommerville (2011), un conjunto de definiciones y
protocolos que permite que dos sistemas se comuniquen entre sí sin
conocer los detalles internos de su implementación. El consumo de un
modelo de Inteligencia Artificial mediante su API permite que el
prototipo envíe la imagen del producto a un servicio externo
especializado y reciba la identificación del artículo, delegando en
dicho servicio la complejidad del procesamiento visual.

La calidad de la respuesta de un modelo multimodal depende en gran
medida de la instrucción que se le proporciona. La \textbf{ingeniería
de prompts} es la práctica de diseñar y refinar dicha instrucción para
orientar el comportamiento del modelo y obtener una respuesta precisa y
en un formato estructurado y procesable. En el presente trabajo, el
prompt se diseña para que el modelo devuelva la identificación del
producto y sus atributos de manera que puedan integrarse directamente
con la lógica de gestión de inventario.

% -------
\subsection{Tecnologías de Desarrollo Web}
% -------

El prototipo se construye sobre tecnologías web de código abierto y
amplia adopción. \textbf{Python} es un lenguaje de programación de alto
nivel, interpretado y de propósito general, con un ecosistema maduro
para el desarrollo de servicios y la integración con APIs de
Inteligencia Artificial; sobre él se emplea el framework
\textbf{FastAPI} para exponer la lógica de negocio como una API REST.
La capa de presentación se desarrolla con \textbf{Next.js}, un
framework para la construcción de interfaces web modernas. La
persistencia se sustenta en \textbf{Supabase}, plataforma basada en el
sistema gestor de base de datos relacional PostgreSQL, que además
provee servicios de sincronización en tiempo real. El despliegue se
realiza en plataformas de capa gratuita, lo que hace viable el proyecto
sin inversión en infraestructura.

% -------
\subsection{Bases de Datos para Gestión de Inventarios}
% -------

Una base de datos, según Date (2001), es un conjunto organizado de
datos relacionados entre sí, almacenados de manera persistente y
accesibles mediante un sistema de gestión (SGBD). En el ámbito de los
sistemas de inventario, la base de datos centraliza toda la información
referente a productos, entradas, salidas, alertas y reportes,
garantizando la integridad y consistencia de los datos. Román Veliz y
Arce Ríos (2023) demuestran en su investigación que el diseño adecuado
de una base de datos relacional para gestión de inventarios es
determinante para la generación de alertas automáticas de stock mínimo
y para la producción de dashboards gerenciales en tiempo real, elementos
que elevan la eficiencia operativa y facilitan la toma de decisiones.

% -------
\subsection{Sistemas de Información Gerencial}
% -------

Los Sistemas de Información Gerencial (SIG) son, según Laudon y Laudon
(2016), sistemas que proporcionan a los gerentes informes periódicos y
acceso en línea al desempeño actual e histórico de la organización. En
el marco del presente proyecto, el sistema propuesto incorpora un módulo
de reportes gerenciales que transforma los datos crudos del inventario
en indicadores clave de gestión (KPIs), como la tasa de rotación de
inventario, el nivel de servicio y el costo de mantenimiento del stock,
facilitando la toma de decisiones estratégicas por parte de la
administración de la empresa.

% -------------------------------------------------------
\section{Bases Legales}
% -------------------------------------------------------

Las bases legales comprenden el conjunto de normas, leyes y
disposiciones jurídicas venezolanas que regulan el desarrollo,
implementación y uso de sistemas de información, tecnología y comercio
electrónico en el país, y que guardan relación directa con el objeto
de estudio.

\textbf{Constitución de la República Bolivariana de Venezuela (1999).}
El Artículo 110 establece que el Estado reconocerá el interés público
de la ciencia, la tecnología, el conocimiento, la innovación y sus
aplicaciones y los servicios de información necesarios, por ser
instrumentos fundamentales para el desarrollo económico, social y
político del país. Este artículo ampara el uso de tecnologías avanzadas
como la Inteligencia Artificial y los servicios digitales en proyectos
de desarrollo nacional y regional.

\textbf{Decreto con Rango, Valor y Fuerza de Ley sobre Mensajes de
Datos y Firmas Electrónicas (2001).} Este decreto regula el uso de los
medios electrónicos y telemáticos en el territorio venezolano, otorgando
plena validez jurídica a los documentos y registros electrónicos. El
Artículo 4 establece que los mensajes de datos tendrán la misma eficacia
probatoria que la ley otorga a los documentos escritos. En el contexto
del sistema propuesto, los registros digitales de inventario generados
tendrán respaldo legal bajo los términos de este decreto.

\textbf{Ley Especial contra los Delitos Informáticos (2001).} Esta ley
tipifica los delitos relacionados con el uso indebido de tecnologías
de la información y establece las sanciones correspondientes. Es
relevante para el proyecto en tanto obliga a garantizar la seguridad,
la confidencialidad y la integridad de los datos almacenados en el
sistema de inventarios, protegiéndolos contra accesos no autorizados,
manipulación ilícita o destrucción de la información.

\textbf{Ley Orgánica de Ciencia, Tecnología e Innovación -- LOCTI
(2010).} En su Artículo 1 establece los lineamientos que orientarán
las políticas y estrategias para la actividad científica, tecnológica
e innovadora, promoviendo el desarrollo nacional de capacidades
tecnológicas propias. El proyecto se enmarca dentro de los principios
de esta ley al proponer una solución tecnológica innovadora desarrollada
con recursos y talento nacional, orientada a resolver una problemática
concreta del sector productivo venezolano.

% -------------------------------------------------------
\section{Referencia Conceptual}
% -------------------------------------------------------

La referencia conceptual o definición de términos básicos recoge los
conceptos clave que, en el contexto de la presente investigación,
adquieren un significado técnico específico. Se presentan en orden
alfabético según las normas institucionales.

\textbf{API (Interfaz de Programación de Aplicaciones):} Conjunto de
definiciones y protocolos que permite que dos sistemas de software se
comuniquen entre sí sin conocer los detalles internos de su
implementación (Sommerville, 2011).

\textbf{Aplicación web:} Programa informático al que se accede a través
de un navegador, cuya lógica se ejecuta en un servidor y cuya interfaz
se presenta en el cliente, sin requerir instalación en el equipo del
usuario (Sommerville, 2011).

\textbf{Código QR:} Código de barras bidimensional que almacena
información ---como un enlace o un identificador de sesión--- y que,
al ser escaneado con la cámara de un dispositivo, permite vincularlo
de forma rápida con un sistema.

\textbf{Control de inventarios:} Conjunto de actividades y
procedimientos destinados a supervisar, registrar y administrar las
existencias de productos de una organización, con el fin de garantizar
su disponibilidad y minimizar los costos de almacenamiento
(Guerrero, 2009).

\textbf{Digitalización:} Proceso de conversión de información, procesos
o registros analógicos a formatos digitales procesables por sistemas
informáticos, con el propósito de automatizar su gestión y facilitar
su acceso y análisis (Laudon y Laudon, 2016).

\textbf{Documento de Requerimientos de Producto (PRD):} Documento que
describe de forma detallada las funcionalidades, restricciones y
condiciones que debe satisfacer un producto o sistema de software,
sirviendo como referencia para su desarrollo y validación (IEEE, 2018).

\textbf{Especificación de Requerimientos de Software (SRS):} Documento
formal que describe las capacidades, restricciones y características
que debe poseer un sistema de software, elaborado siguiendo los
lineamientos del estándar IEEE 830 (IEEE, 2018).

\textbf{Indicadores Clave de Gestión (KPIs):} Métricas cuantitativas y
cualitativas utilizadas para evaluar el desempeño de los procesos de una
organización en relación con sus objetivos estratégicos (Laudon y
Laudon, 2016).

\textbf{Ingeniería de prompts:} Práctica de diseñar y refinar la
instrucción (prompt) que se proporciona a un modelo de Inteligencia
Artificial para orientar su comportamiento y obtener una respuesta
precisa y en un formato procesable.

\textbf{Inteligencia Artificial (IA):} Disciplina de la computación
orientada al diseño de sistemas capaces de simular capacidades cognitivas
humanas como el aprendizaje, el razonamiento y el reconocimiento de
patrones (Russell y Norvig, 2016).

\textbf{Inventario:} Conjunto de bienes tangibles que una empresa
mantiene en existencia para ser utilizados en la producción o para su
venta directa, incluyendo materias primas, productos en proceso y
productos terminados (Ballou, 2004).

\textbf{Modelo de lenguaje multimodal:} Modelo de Inteligencia
Artificial capaz de recibir y relacionar simultáneamente distintas
modalidades de información ---como texto e imagen--- y producir una
respuesta en lenguaje natural, permitiendo identificar el contenido de
una imagen a partir de una instrucción textual (Russell y Norvig,
2016).

\textbf{Reconocimiento de imágenes:} Capacidad de un sistema
computacional para identificar y clasificar el contenido de una imagen
digital, asignando etiquetas o categorías a los elementos visuales
presentes en ella (Russell y Norvig, 2016).

\textbf{Sincronización en tiempo real:} Mecanismo mediante el cual un
cambio efectuado en un dispositivo se refleja de manera inmediata en
los demás dispositivos conectados al mismo sistema, sin necesidad de
recargar manualmente la información.

\textbf{Sistema de información:} Conjunto integrado de componentes para
la recolección, almacenamiento, procesamiento y comunicación de datos,
con el fin de apoyar la toma de decisiones y el control organizacional
(Laudon y Laudon, 2016).

\textbf{Stock:} Cantidad de mercancías o productos que una empresa
mantiene almacenados para hacer frente a la demanda futura o a las
necesidades del proceso productivo (Ballou, 2004).

% ============================================================


% -------------------------------------------------------
%  CAPÍTULO III  –  MARCO METODOLÓGICO
% -------------------------------------------------------
\chapter{Marco Metodológico}

El marco metodológico es el apartado del informe de investigación donde
se expone la manera como se realiza el estudio, los pasos seguidos para
llevarlo a cabo y los instrumentos empleados para responder al problema
planteado. En este capítulo se desarrollan el tipo y el diseño de la
investigación, la población y la muestra, las técnicas e instrumentos de
recolección de datos, la operacionalización de la variable, el
procedimiento para el logro de los objetivos y las técnicas de
procesamiento y análisis de los datos.

% -------------------------------------------------------
\section{Tipo de Investigación}
% -------------------------------------------------------

De acuerdo con el nivel de profundidad con el que se aborda el objeto de
estudio, la presente investigación se clasifica como de tipo
\textbf{Descriptiva}. Según Arias (2012), la investigación descriptiva
consiste en la caracterización de un hecho, fenómeno, individuo o grupo,
con el fin de establecer su estructura o comportamiento. En este sentido,
el estudio es descriptivo porque, en su fase inicial, caracteriza y
diagnostica la situación actual del proceso de registro y control de
inventarios del Bodegón Pa' Donde Natty C.A., identificando las brechas
de trazabilidad, los errores recurrentes y los requerimientos
funcionales y técnicos que debe satisfacer el prototipo propuesto.

Por sus propósitos, la investigación se enmarca en la modalidad
\textbf{Aplicada}, dado que persigue generar un producto tecnológico
concreto ---un prototipo funcional de un sistema de digitalización y
registros de productos con Inteligencia Artificial--- orientado a la
solución directa de un problema real identificado en la empresa.
Asimismo, es de carácter \textbf{Proyectivo}, puesto que, conforme al
marco de los proyectos propositivos definido por Hurtado (2010),
propone, desarrolla y valida una solución tecnológica como alternativa
al problema de gestión de inventarios detectado.

% -------------------------------------------------------
\section{Diseño de la Investigación}
% -------------------------------------------------------

En atención a la estrategia general adoptada para responder al problema
planteado, la investigación se enmarca en un diseño \textbf{de Campo}.
Según Arias (2020), la investigación de campo es aquella que consiste en
la recolección de datos directamente de los sujetos investigados o de la
realidad donde ocurren los hechos, sin manipular o controlar variable
alguna. En este proyecto, el diseño es de campo porque los datos
primarios relativos al proceso actual de inventario, las necesidades del
personal y la percepción de la problemática se obtienen directamente en
las instalaciones del Bodegón Pa' Donde Natty C.A., ubicado en Ciudad
Bolívar, estado Bolívar.

El diseño de campo se complementa con un componente \textbf{Documental},
dado que parte del estudio requiere la búsqueda, recuperación, análisis
e interpretación de fuentes bibliográficas, hemerográficas y
electrónicas relacionadas con la Inteligencia Artificial, los modelos
de lenguaje multimodales, el consumo de servicios mediante API, Python
y los sistemas de gestión de inventarios.
Arias (2017) define la investigación documental como aquella basada en
la obtención y análisis de datos provenientes de materiales impresos u
otros tipos de documentos, los cuales sustentan teórica y técnicamente
la construcción del prototipo.

Finalmente, la investigación es de tipo \textbf{no experimental}, en
tanto que el investigador no somete el objeto de estudio a condiciones
de experimentación controlada sobre variables poblacionales, sino que
observa y analiza la situación tal como se presenta para, posteriormente,
desarrollar y validar la solución propuesta en un entorno controlado de
desarrollo.

% -------------------------------------------------------
\section{Población y Muestra}
% -------------------------------------------------------

\subsection{Población}

Según Arias (2012), la población es cualquier conjunto finito o infinito
de elementos con características comunes para los cuales serán extensivas
las conclusiones de la investigación. En el contexto del presente
estudio, la población está constituida por la totalidad del personal que
labora en el Bodegón Pa' Donde Natty C.A. y que participa directa o
indirectamente en los procesos de registro, almacenamiento y control de
inventarios, conformada por cinco (05) trabajadores: el gerente general,
el encargado de almacén, dos (02) operarios de recepción y despacho y el
personal administrativo responsable del registro de entradas y salidas
de productos.

\subsection{Muestra}

Dado que la población de estudio es finita, conocida y de tamaño
reducido, se aplica el criterio de \textbf{censo poblacional}; es decir,
se trabaja con la totalidad de los sujetos que la conforman, sin recurrir
a técnicas de muestreo probabilístico. Según Tamayo y Tamayo (2007),
cuando la población es pequeña y accesible en su totalidad, lo más
conveniente es incorporar a todos sus integrantes como unidad de
análisis, a fin de garantizar la máxima representatividad de los datos
recolectados. En consecuencia, la muestra está conformada por los mismos
cinco (05) trabajadores que integran la población.

\subsection{Unidad de Análisis}

La unidad de análisis es el Bodegón Pa' Donde Natty C.A.,
específicamente su área de almacén e inventarios ubicada en Ciudad
Bolívar, estado Bolívar. Según Arias (2006), la unidad de investigación
es por lo general la unidad organizativa en donde tiene lugar el
problema, pudiendo hacerse extensiva a cualquier situación a mejorar
dentro de la realidad organizacional.

% -------------------------------------------------------
\section{Técnicas e Instrumentos de Recolección de Datos}
% -------------------------------------------------------

Las técnicas de recolección de datos son, según Arias (2012), las
distintas formas o maneras de obtener la información, mientras que los
instrumentos son los medios materiales que se emplean para recoger y
almacenar dicha información. A continuación se describen las técnicas
seleccionadas y el instrumento asociado a cada una, indicando el
objetivo específico al que sirven.

\subsection{Técnica: La Encuesta}

La encuesta es, según Arias (2012), una técnica que pretende obtener
información que suministra un grupo o muestra de sujetos acerca de sí
mismos o en relación con un tema en particular. Se emplea para recabar
de manera sistemática y estandarizada las opiniones, percepciones y
necesidades del personal del Bodegón Pa' Donde Natty C.A. en relación
con el proceso de registro y control de inventarios. Sirve,
principalmente, al \textbf{Objetivo Específico N° 1}.

\noindent\textbf{Instrumento: Cuestionario.} Según Arias (2012,
p.\,72), el cuestionario es la modalidad de encuesta que se realiza de
forma escrita mediante un formato contentivo de una serie de preguntas.
El cuestionario aplicado es estructurado, con preguntas cerradas y
escala de frecuencia tipo Likert, diseñado bajo los indicadores del
cuadro de operacionalización para medir la frecuencia de errores en el
registro de productos, el tiempo empleado en los conteos de inventario y
las dificultades más recurrentes del proceso actual.

\subsection{Técnica: Observación Directa}

Según Tamayo (2007, p.\,24), la observación directa es aquella en la
cual el investigador puede observar y recoger datos mediante su propia
observación, sin intermediación alguna. Se aplica en las instalaciones
de la empresa para apreciar el proceso actual de registro, el estado del
área de almacén y el flujo de entradas y salidas de mercancía. Sirve al
\textbf{Objetivo Específico N° 1}, y se complementa con una entrevista
no estructurada al gerente general y al encargado de almacén orientada
a indagar las debilidades del sistema actual y las expectativas respecto
al prototipo.

\noindent\textbf{Instrumentos: Guía de Observación Estructurada y Guía
de Entrevista.} La guía de observación especifica de antemano los
aspectos a registrar (método de conteo, herramientas de registro,
disposición física de los productos y frecuencia de actualización del
inventario); la guía de entrevista contiene los tópicos a abordar,
apoyada en una libreta de notas.

\subsection{Técnica: Análisis Documental}

El análisis documental es definido por Arias (2012, p.\,27) como un
proceso basado en la búsqueda, recuperación, análisis, crítica e
interpretación de datos secundarios registrados en fuentes
documentales. En el contexto de la ingeniería de software, esta técnica
se aplica de forma sistemática sobre el documento de Especificación de
Requerimientos de Software (SRS) para extraer, categorizar y validar las
reglas de negocio, las restricciones técnicas y los requerimientos
previamente diagnosticados. Sirve, principalmente, al \textbf{Objetivo
Específico N° 2}.

\noindent\textbf{Instrumento: Matriz de Trazabilidad de Requisitos
(MTR).} De acuerdo con Sommerville (2011), la trazabilidad de
requerimientos consiste en registrar explícitamente las relaciones de
correspondencia entre cada requisito del sistema, su origen y los
componentes de diseño y de código que lo implementan. La MTR se
estructura como una tabla de doble entrada cuyas filas corresponden a
cada requerimiento funcional extraído de la SRS y cuyas columnas
verifican: el componente arquitectónico responsable, la entidad o tabla
de la base de datos relacional asociada y la pantalla o vista de la
interfaz de usuario que lo materializa.

\subsection{Técnica: Observación Sistemática y Pruebas de Software}

La técnica seleccionada para evaluar el alcance del cuarto objetivo
específico es la observación sistemática, ejecutada mediante pruebas
funcionales aplicadas al prototipo. Según Arias (2012), la observación
sistemática permite recopilar datos de forma estructurada y con
objetivos previamente determinados. En ingeniería de software,
Pressman y Maxim (2020) señalan que las pruebas funcionales en
prototipos consisten en operar el software en un entorno controlado
para evaluar si responde fielmente a los estímulos de entrada según las
especificaciones del diseño, sin que ello implique su implantación
definitiva en el escenario real. En este estudio, las pruebas se
ejecutan en un entorno local de desarrollo, observando y verificando
matemáticamente el comportamiento de los algoritmos de reconocimiento de
productos y de gestión de stock. Sirve, principalmente, al
\textbf{Objetivo Específico N° 4}.

\noindent\textbf{Instrumento: Lista de Cotejo (Checklist).} Arias
(2012) define la lista de cotejo como un instrumento que registra la
presencia o ausencia de características específicas mediante una escala
dicotómica (Cumple / No Cumple). Se estructura como un mecanismo de
control de calidad técnica cuyos ítems verifican, de forma estática y
dinámica, la renderización correcta de las interfaces, la persistencia
de los datos en la base de datos y la exactitud del algoritmo de
reconocimiento y registro ante casos de prueba previamente definidos.

\subsection{Metodología de Desarrollo y Entorno de Trabajo}

El desarrollo del prototipo se gestiona mediante la metodología ágil
\textbf{Kanban}, cuyo flujo de trabajo se visualiza en un tablero
dividido en columnas (Por Hacer, En Desarrollo, En Pruebas, Finalizado),
lo que permite controlar de manera incremental el avance de cada
requerimiento y tarea técnica. El prototipo se construye como una
\textbf{aplicación web} con frontend y backend desacoplados. La
\textbf{aplicación principal} se ejecuta en el navegador de una
\textbf{computadora (PC)} ---desarrollada con \textbf{Next.js}--- y
concentra la gestión del inventario: panel de indicadores, registro,
control de existencias, alertas y reportes. El \textbf{teléfono móvil
actúa únicamente como dispositivo de captura}: mediante un código QR
generado por la aplicación de la PC, el teléfono se vincula a la sesión
activa y abre una vista ligera cuya única función es tomar la fotografía
del producto con su cámara y enviarla al sistema. La imagen capturada
desde el teléfono se transmite al backend y el resultado se refleja en
tiempo real en la aplicación de la PC. La capa de lógica de negocio se
implementa en \textbf{Python} con el framework \textbf{FastAPI}, que
expone una API REST; y la capa de persistencia se sustenta en
\textbf{Supabase}, plataforma basada en el sistema gestor de base de
datos relacional \textbf{PostgreSQL}, que además provee la
sincronización en tiempo real entre el teléfono y la PC. El
reconocimiento de los productos no se realiza con un modelo entrenado
localmente, sino mediante la \textbf{interfaz de programación de
aplicaciones (API) de un modelo de lenguaje multimodal} de Inteligencia
Artificial, lo que traslada la complejidad del procesamiento visual a un
servicio especializado y permite centrar el aporte del trabajo en la
construcción del sistema de gestión de inventario. El despliegue se
realiza en plataformas de capa gratuita: el frontend en \textbf{Vercel}
y el backend en \textbf{Render}. Como instrumentos de apoyo se emplean
un teléfono con cámara para la captura de imágenes, una computadora para
el desarrollo, una libreta de notas y unidades de almacenaje para el
respaldo de los datos y archivos.

\subsection{Integración del Servicio de Reconocimiento (API
Multimodal)}

A diferencia de un sistema de información convencional, el presente
prototipo delega el reconocimiento de productos en un modelo de lenguaje
multimodal de Inteligencia Artificial, consumido a través de su API en
su capa gratuita. La metodología para incorporar dicho servicio
contempla las siguientes fases:

\begin{enumerate}
  \item \textbf{Selección y registro del servicio:} se selecciona un
        proveedor de modelo multimodal con capa gratuita y soporte de
        comprensión de imágenes, y se obtiene la credencial de acceso
        (clave de API).

  \item \textbf{Diseño de la instrucción (prompt):} se elabora la
        instrucción que se envía junto con la imagen del producto, de
        modo que el modelo devuelva, en un formato estructurado, la
        identificación del producto y sus atributos relevantes.

  \item \textbf{Implementación del cliente de la API:} se desarrolla, en
        la capa de lógica de negocio (Python/FastAPI), el componente que
        envía la imagen capturada al servicio, recibe la respuesta y la
        interpreta.

  \item \textbf{Acoplamiento con la gestión de stock:} el producto
        reconocido se asocia con el catálogo y se registra el
        movimiento de inventario en la base de datos PostgreSQL.

  \item \textbf{Evaluación del reconocimiento:} se mide la exactitud del
        servicio sobre un conjunto de imágenes de prueba de los
        productos del bodegón, calculando el porcentaje de
        reconocimientos correctos y el tiempo de respuesta.
\end{enumerate}

% -------------------------------------------------------
\section{Operacionalización de la Variable}
% -------------------------------------------------------

La operacionalización es el proceso mediante el cual una variable
compleja se descompone en conceptos concretos, observables y medibles,
es decir, en dimensiones e indicadores. La \textbf{variable de estudio}
es el \textit{prototipo de un sistema de digitalización y registros de
productos con Inteligencia Artificial para el control de inventarios},
entendida como el conjunto de necesidades, componentes de diseño,
módulos desarrollados y criterios de cumplimiento que permiten
automatizar el registro y el control de existencias de la empresa
mediante visión computacional e Inteligencia Artificial. A continuación
se presenta el Cuadro de Operacionalización de la Variable.

\vspace{12pt}

% Tabla de operacionalización
\setstretch{1.0}
\scriptsize
\hbadness=10000
\setlength{\LTleft}{0pt}
\setlength{\LTright}{0pt}
\setlength{\tabcolsep}{2pt}
\begin{longtable}{|>{\raggedright\arraybackslash}p{2.5cm}
                  |>{\raggedright\arraybackslash}p{2.1cm}
                  |>{\raggedright\arraybackslash}p{4.0cm}
                  |>{\raggedright\arraybackslash}p{2.9cm}|}
\caption{Operacionalización de la Variable.}\\
\hline
\textbf{Objetivo Específico} &
\textbf{Dimensión} &
\textbf{Indicadores} &
\textbf{Técnica / Instrumento} \\
\hline
\endfirsthead
\hline
\textbf{Objetivo Específico} &
\textbf{Dimensión} &
\textbf{Indicadores} &
\textbf{Técnica / Instrumento} \\
\hline
\endhead
\hline
\endfoot

1. Diagnosticar las necesidades funcionales y técnicas del proceso
actual de registro y control de inventarios. &
Operativa &
-- Tipo y frecuencia de errores de registro \newline
-- Tiempo promedio por ciclo de inventario \newline
-- Necesidades tecnológicas priorizadas &
Encuesta (cuestionario); Observación directa; Entrevista \\
\hline

2. Diseñar la arquitectura lógica, el modelo de datos y la interfaz de
usuario (SRS). &
Técnica (diseño) &
-- Índice de acoplamiento de la arquitectura \newline
-- Nivel de normalización de la base de datos (3FN) \newline
-- Porcentaje de cobertura de interfaces respecto a la SRS &
Análisis documental; Matriz de Trazabilidad de Requisitos (MTR) \\
\hline

3. Desarrollar el prototipo funcional integrando el reconocimiento
mediante la API de un modelo multimodal. &
Funcional (desarrollo) &
-- Porcentaje de módulos desarrollados \newline
-- Integración correcta del servicio de reconocimiento (API) \newline
-- Efectividad del reconocimiento y del registro &
Metodología Kanban; Integración de la API; Observación
sistemática \\
\hline

4. Validar el funcionamiento del prototipo mediante pruebas en un
entorno controlado. &
Validación &
-- Tasa de cumplimiento de requisitos funcionales \newline
-- Exactitud del reconocimiento (igual o superior al 90\,\%) \newline
-- Tiempo de respuesta (menor o igual a 3 s) &
Observación sistemática (pruebas de caja negra); Lista de cotejo \\
\hline
\end{longtable}
\hbadness=1000
\normalsize
\setlength{\tabcolsep}{6pt}
\setstretch{1.5}

\noindent\small\textit{Fuente: Elaboración propia (2026).}
\normalsize

% -------------------------------------------------------
\section{Procedimiento para el Logro de los Objetivos}
% -------------------------------------------------------

En esta sección se describe de manera detallada y secuencial el conjunto
de actividades que se ejecutan para alcanzar cada uno de los objetivos
específicos planteados en el Capítulo I. Para cada actividad se indica
entre paréntesis el indicador del cuadro de operacionalización al que
responde, garantizando la trazabilidad entre la metodología y los
resultados presentados en el Capítulo IV.

\noindent\textbf{Objetivo Específico N° 1.} Diagnosticar las
necesidades funcionales y técnicas del proceso actual de registro y
control de inventarios.

\begin{enumerate}
  \item \textbf{Diseño de los instrumentos:} se elaboran el
        cuestionario, la guía de observación y la guía de entrevista,
        derivados de los indicadores operativos de la variable.
        \textit{(Indicadores: tipo y frecuencia de errores de registro;
        necesidades tecnológicas priorizadas.)}

  \item \textbf{Aplicación de los instrumentos:} se realizan visitas
        técnicas a la empresa para aplicar la observación directa, las
        entrevistas al gerente general y al encargado de almacén, y el
        cuestionario a los operarios y al personal administrativo.
        \textit{(Indicador: tiempo promedio por ciclo de inventario.)}

  \item \textbf{Tabulación y análisis:} los datos se exportan a una
        hoja de cálculo para generar estadísticas descriptivas que
        permitan identificar los cuellos de botella del proceso actual.
        \textit{(Indicadores: tipo y frecuencia de errores de registro;
        tiempo promedio por ciclo de inventario.)}

  \item \textbf{Especificación de requisitos:} como actividad de cierre,
        se consolidan las necesidades en requerimientos funcionales y
        restricciones técnicas. \textit{(Indicador: necesidades
        tecnológicas priorizadas.)}
\end{enumerate}

\noindent\textbf{Objetivo Específico N° 2.} Diseñar la arquitectura
lógica, el modelo de datos y la interfaz de usuario, consolidados en la
SRS.

\begin{enumerate}
  \item \textbf{Diseño de la arquitectura lógica:} a partir de los
        requerimientos se modela la estructura del sistema bajo un
        patrón modular, buscando la independencia de los componentes.
        \textit{(Indicador: índice de acoplamiento de la arquitectura.)}

  \item \textbf{Diseño del modelo de datos:} se estructura el esquema de
        la base de datos relacional, definiendo entidades, claves y
        restricciones adecuadas para la gestión de inventarios.
        \textit{(Indicador: nivel de normalización de la base de datos
        en Tercera Forma Normal.)}

  \item \textbf{Prototipado de la interfaz de usuario:} se diseñan las
        vistas conceptuales (wireframes) orientadas a un usuario sin
        formación técnica avanzada. \textit{(Indicador: porcentaje de
        cobertura de interfaces respecto a la SRS.)}

  \item \textbf{Mapeo de control:} los productos de diseño se integran
        en la Matriz de Trazabilidad de Requisitos, asegurando la
        correspondencia entre cada requerimiento y los componentes
        diseñados. \textit{(Indicadores: índice de acoplamiento;
        porcentaje de cobertura de interfaces.)}
\end{enumerate}

\noindent\textbf{Objetivo Específico N° 3.} Desarrollar el prototipo
funcional integrando el reconocimiento mediante la API de un modelo
multimodal.

\begin{enumerate}
  \item \textbf{Configuración del entorno:} instalación y despliegue del
        entorno de desarrollo con Python, FastAPI y Supabase como gestor de base de
        datos PostgreSQL. \textit{(Indicador: porcentaje de módulos
        desarrollados.)}

  \item \textbf{Integración del servicio de reconocimiento:} registro
        en el proveedor del modelo multimodal, obtención de la clave de
        API, diseño de la instrucción (prompt) e implementación del
        cliente que envía la imagen y procesa la respuesta.
        \textit{(Indicador: integración correcta del servicio de
        reconocimiento.)}

  \item \textbf{Codificación de los módulos:} desarrollo del módulo de
        captura móvil vinculada por QR (toma de la foto en el
        teléfono), el módulo de reconocimiento (consumo de la API), el módulo de gestión de
        stock y base de datos (PostgreSQL) y el módulo de reportes
        gerenciales, gestionados como tarjetas del tablero Kanban a
        partir de la MTR. \textit{(Indicador: porcentaje de módulos
        desarrollados.)}

  \item \textbf{Integración:} acoplamiento del servicio de
        reconocimiento con la aplicación web, programando el registro
        automatizado y la actualización del inventario en la base de
        datos. \textit{(Indicador: efectividad del reconocimiento y del
        registro.)}
\end{enumerate}

\noindent\textbf{Objetivo Específico N° 4.} Validar el funcionamiento
del prototipo mediante pruebas en un entorno controlado.

\begin{enumerate}
  \item \textbf{Diseño de los casos de prueba:} se definen los casos de
        prueba de caja negra a partir de los requerimientos funcionales
        de la SRS. \textit{(Indicador: tasa de cumplimiento de
        requisitos funcionales.)}

  \item \textbf{Ejecución de las pruebas:} se opera el prototipo en el
        entorno controlado y se registran las evidencias en la lista de
        cotejo, evaluando el cumplimiento de cada requerimiento.
        \textit{(Indicadores: tasa de cumplimiento de requisitos
        funcionales; tiempo de respuesta.)}

  \item \textbf{Verificación de la exactitud del reconocimiento:} se
        somete al servicio un conjunto de imágenes de prueba de los
        productos del bodegón y se calcula el porcentaje de
        reconocimientos correctos, contrastándolo con la identificación
        verificada manualmente. \textit{(Indicador: exactitud del
        reconocimiento.)}
\end{enumerate}

% -------------------------------------------------------
\section{Técnicas de Procesamiento y Análisis de los Datos}
% -------------------------------------------------------

Considerando la naturaleza mixta del estudio, el procesamiento y
análisis de los datos se organiza en dos vertientes: la estadística
descriptiva aplicada a los datos de la fase de diagnóstico, y el
análisis de cobertura y conformidad aplicado a los instrumentos
técnicos de control (MTR y lista de cotejo). Según Arias (2006), las
técnicas de análisis son las técnicas lógicas (inducción, deducción,
análisis-síntesis) o estadísticas (descriptivas o inferenciales) que se
emplean para descifrar lo que revelan los datos recolectados.

\subsection{Procesamiento y Análisis del Diagnóstico (Objetivo N° 1)}

Los datos primarios recopilados mediante el cuestionario y la
observación se someten a técnicas de \textbf{estadística descriptiva}.
Se calculan frecuencias absolutas y porcentajes relativos para cada
indicador evaluado, y los resultados se representan mediante tablas y
gráficos estadísticos de barras o sectores. La información cualitativa
proveniente de las entrevistas y la observación se interpreta de manera
contextual, relacionándola con los fundamentos teóricos del Capítulo II,
con el fin de transformar los hallazgos en las pautas funcionales de la
ingeniería de requerimientos del prototipo.

\subsection{Análisis del Diseño (Objetivo N° 2)}

Sobre los productos de diseño se aplica un \textbf{análisis de cobertura
lógica} mediante la Matriz de Trazabilidad de Requisitos: se verifica
que cada requerimiento de la SRS posea un componente arquitectónico, una
entidad de base de datos y una pantalla de interfaz asignados,
garantizando un diseño cohesionado, con bajo índice de acoplamiento y
sin requerimientos huérfanos. Asimismo, se verifica que el modelo de
datos cumpla con la Tercera Forma Normal.

\subsection{Análisis del Desarrollo y la Validación (Objetivos N° 3 y
N° 4)}

El avance de la construcción se procesa calculando el
\textbf{porcentaje de módulos desarrollados} respecto al total
planificado. La exactitud del reconocimiento se evalúa sometiendo al
servicio multimodal un conjunto de imágenes de prueba de los productos
del bodegón y calculando el \textbf{porcentaje de reconocimientos
correctos} frente a la identificación verificada manualmente. Para la
validación funcional, los datos dicotómicos de la lista de cotejo
(Cumple / No Cumple) se tabulan para determinar la \textbf{tasa de
cumplimiento de requisitos funcionales}, expresada como la proporción
porcentual de ítems aprobados frente al total exigido. El criterio de
aceptación exige una exactitud de reconocimiento igual o superior al
90\,\% y un margen de error nulo en el registro automatizado de las
existencias.

% -------------------------------------------------------
\section{Cronograma de Actividades}
% -------------------------------------------------------

El desarrollo del presente Trabajo Especial de Grado se planificó en un
horizonte temporal de cinco (05) meses, comprendido entre enero y mayo de
2026, distribuyendo las actividades conforme a la secuencia de los cuatro
objetivos específicos. El cronograma que se presenta a continuación
organiza dichas actividades en el tiempo; cada barra señala el período de
ejecución de la actividad correspondiente.

\vspace{12pt}

\newcommand{\barragantt}{\rule{0.7cm}{6pt}}
\setstretch{1.0}
\scriptsize
\setlength{\tabcolsep}{3pt}
\begin{longtable}{|>{\raggedright\arraybackslash}p{6.3cm}|c|c|c|c|c|}
\caption{Cronograma de actividades del Trabajo Especial de Grado.}\\
\hline
\textbf{Actividad} & \textbf{Ene} & \textbf{Feb} & \textbf{Mar} &
\textbf{Abr} & \textbf{May} \\
\hline
\endfirsthead
\hline
\textbf{Actividad} & \textbf{Ene} & \textbf{Feb} & \textbf{Mar} &
\textbf{Abr} & \textbf{May} \\
\hline
\endhead
\hline
\endfoot
Revisión documental, antecedentes y fundamentos teóricos &
\barragantt & \barragantt & & & \\
\hline
Diseño y aplicación de los instrumentos de diagnóstico (Obj.\ 1) &
& \barragantt & \barragantt & & \\
\hline
Análisis del diagnóstico y especificación de requerimientos (Obj.\ 1) &
& & \barragantt & & \\
\hline
Diseño de la arquitectura, el modelo de datos y la SRS (Obj.\ 2) &
& & \barragantt & \barragantt & \\
\hline
Desarrollo de los módulos e integración de la API multimodal (Obj.\ 3) &
& & & \barragantt & \barragantt \\
\hline
Pruebas de caja negra y validación del prototipo (Obj.\ 4) &
& & & & \barragantt \\
\hline
Redacción, revisión y corrección del informe final &
\barragantt & \barragantt & \barragantt & \barragantt & \barragantt \\
\hline
\end{longtable}
\setlength{\tabcolsep}{6pt}
\normalsize

\vspace{2pt}
\noindent\small\textit{Fuente: Elaboración propia (2026).}\normalsize
\setstretch{1.5}

% -------------------------------------------------------
%  CAPÍTULO IV  –  ANÁLISIS DE RESULTADOS
% -------------------------------------------------------
\chapter{Análisis de Resultados}

En el presente capítulo se exponen, analizan e interpretan los
resultados obtenidos durante el desarrollo de la investigación. No se
trata de una mera presentación de tablas y figuras: cada resultado se
discute a la luz de los indicadores del Cuadro de Operacionalización
de la Variable formulado en el Capítulo III y se conecta con los
fundamentos teóricos del Capítulo II, de modo que el lector pueda
comprender no solo qué se obtuvo, sino qué significa cada hallazgo y
qué implicaciones tiene para el problema planteado en el Capítulo I.
La exposición sigue la secuencia lógica de los cuatro objetivos
específicos ---diagnóstico, diseño, desarrollo y validación---; cada
objetivo se introduce indicando los instrumentos e indicadores con que
se abordó, presenta sus tablas y figuras precedidas de un llamado en el
texto y seguidas de su análisis interpretativo, y cierra con un párrafo
que afirma y justifica su cumplimiento con base en la evidencia. Al
final del capítulo se incorpora, de forma complementaria, un breve
análisis de factibilidad del prototipo.

% =======================================================
\section{Diagnóstico de las Necesidades Funcionales y Técnicas}
% =======================================================

Este apartado responde al \textbf{primer objetivo específico}:
diagnosticar las necesidades funcionales y técnicas del proceso actual
de registro y control de inventarios del Bodegón Pa' Donde Natty C.A.
Conforme a la dimensión operativa del Cuadro de Operacionalización, el
diagnóstico se apoyó en tres instrumentos complementarios ---la guía
de observación estructurada, la entrevista no estructurada al gerente
general y al encargado de almacén, y el cuestionario aplicado al censo
poblacional de cinco (05) trabajadores--- y se organizó según sus tres
indicadores: el tipo y la frecuencia de errores de registro, el tiempo
promedio por ciclo de inventario y las necesidades tecnológicas
priorizadas. La triangulación de estas tres fuentes busca que las
conclusiones del diagnóstico no descansen en un único instrumento, sino
en la convergencia de la observación del investigador, la voz directiva
y la percepción del personal operativo.

\subsection{Proceso de Inventario Actual}

A través de la observación directa realizada en las instalaciones de
la empresa se constató que el control de existencias se efectúa de
forma manual, mediante conteos físicos periódicos, anotaciones en
papel y el uso parcial de hojas de cálculo no integradas entre sí. No
existe un registro centralizado ni una actualización en tiempo real de
los movimientos de mercancía. Esta constatación es relevante porque
confirma de forma empírica la problemática descrita en el Capítulo I y,
al mismo tiempo, ilustra con un caso concreto lo que Laudon y Laudon
(2016) describen como un proceso analógico: un flujo de trabajo que
todavía no ha sido digitalizado y que, por ello, arrastra error humano,
lentitud y ausencia de trazabilidad. El hecho de que coexistan papel y
hojas de cálculo no integradas no constituye una digitalización
parcial beneficiosa, sino una fuente adicional de discrepancia, pues la
información se fragmenta en soportes que no se sincronizan entre sí.

En la Tabla 1 se presenta la frecuencia con que el personal encuestado
reporta la ocurrencia de errores en el proceso actual de inventario,
obtenida a partir del cuestionario aplicado al censo poblacional de
cinco (05) trabajadores.

\vspace{6pt}
\setstretch{1.0}
\noindent\textbf{Tabla 1}\\
\textit{Frecuencia de errores reportados en el proceso de inventario
actual}
\vspace{4pt}

\begin{center}
\begin{tabular}{|>{\raggedright\arraybackslash}p{6.0cm}|c|c|}
\hline
\textbf{Tipo de error} & \textbf{Frecuencia (f)} &
\textbf{Porcentaje (\%)} \\
\hline
Discrepancia entre inventario físico y registrado & 5 & 100 \\
\hline
Desabastecimiento de productos de alta rotación & 4 & 80 \\
\hline
Pérdida de mercancía por vencimiento no detectado & 4 & 80 \\
\hline
Duplicidad de registros al transcribir datos & 3 & 60 \\
\hline
Retardo en la reposición de stock & 3 & 60 \\
\hline
\end{tabular}
\end{center}
\vspace{2pt}
\noindent\small\textit{Fuente: Elaboración propia (2026), con base en
el cuestionario aplicado.}\normalsize
\setstretch{1.5}

\vspace{6pt}

Los datos de la Tabla 1 admiten una lectura que va más allá de la
simple enumeración de frecuencias. El hecho de que el 100\,\% de los
trabajadores coincida en señalar la discrepancia entre el inventario
físico y el registrado como el error más recurrente no es un dato
aislado: es el síntoma central del problema y la consecuencia directa
de la ausencia de un registro único y en tiempo real. En términos de
Guerrero (2009), el control de inventarios solo cumple su propósito si
garantiza la disponibilidad del producto en el momento y lugar
requeridos; una discrepancia universalmente percibida indica
precisamente que ese propósito no se está cumpliendo. Los dos errores
que le siguen ---desabastecimiento de productos de alta rotación y
pérdidas por vencimiento, ambos con 80\,\%--- son, a su vez, efectos
encadenados de la misma causa raíz: sin información actualizada, el
personal no puede anticipar ni la rotura ni la caducidad. Esta lectura
es importante para el diseño posterior, porque revela que no se
requieren funciones aisladas, sino un sistema que ataque la causa común
(la falta de un registro centralizado y oportuno) de la que se derivan
los demás errores. Los dos errores de menor frecuencia ---duplicidad de
registros y retardo en la reposición, con 60\,\%--- confirman además
que la transcripción manual introduce ruido en los datos, justificando
la necesidad de automatizar la captura.

En cuanto al tiempo invertido en el ciclo de inventario, la Figura 1
resume la distribución porcentual del tiempo promedio reportado por el
personal para completar un ciclo de conteo físico.

\figuratesis{uploads/fig1_tiempos.pdf}{0.78}{Figura 1. Distribución del
tiempo promedio por ciclo de inventario manual.}

\noindent\small\textit{Fuente: Elaboración propia (2026), con base en
el cuestionario aplicado.}\normalsize

\vspace{4pt}

La interpretación de la Figura 1 conecta directamente con el indicador
``tiempo promedio por ciclo de inventario'' de la dimensión operativa.
El ciclo completo de inventario manual requiere en promedio entre
cuatro (04) y seis (06) horas de trabajo, durante las cuales la
operación regular del negocio se ve parcialmente interrumpida. Más allá
de la cifra, lo significativo es su implicación económica: en un
negocio de alta rotación como el estudiado, cada ciclo de conteo no es
solo tiempo de personal, sino tiempo de operación degradada. Este
hallazgo aporta una justificación cuantitativa al objetivo del sistema,
pues cualquier mecanismo que reduzca la duración del conteo ---como el
reconocimiento automatizado del producto que se propone--- se traduce de
forma inmediata en ahorro operativo, y no solo en una mejora de
exactitud.

\subsection{Necesidades Tecnológicas Identificadas}

A partir de las entrevistas no estructuradas sostenidas con el gerente
general y el encargado de almacén se identificaron las necesidades
tecnológicas priorizadas que el personal directivo espera satisfacer
con el prototipo. En la Tabla 2 se sintetizan dichas necesidades,
ordenadas según la prioridad asignada por la dirección.

\vspace{6pt}
\setstretch{1.0}
\noindent\textbf{Tabla 2}\\
\textit{Necesidades tecnológicas identificadas, según prioridad}
\vspace{4pt}

\begin{center}
\begin{tabular}{|c|>{\raggedright\arraybackslash}p{8.0cm}|c|}
\hline
\textbf{N°} & \textbf{Necesidad tecnológica} & \textbf{Prioridad} \\
\hline
1 & Acceso controlado al sistema mediante autenticación de usuarios &
Alta \\
\hline
2 & Registro digital centralizado y en tiempo real de productos & Alta \\
\hline
3 & Reconocimiento automatizado de productos para agilizar el conteo &
Alta \\
\hline
4 & Control de entradas, salidas y existencias actualizado al
instante & Alta \\
\hline
5 & Alertas automáticas de stock mínimo y de productos por vencer &
Alta \\
\hline
6 & Organización del catálogo por categorías & Media \\
\hline
7 & Reportes gerenciales para la toma de decisiones & Media \\
\hline
8 & Respaldo y configuración del sistema & Media \\
\hline
9 & Reducción de la dependencia del registro manual en papel & Media \\
\hline
\end{tabular}
\end{center}
\vspace{2pt}
\noindent\small\textit{Fuente: Elaboración propia (2026), con base en
las entrevistas aplicadas.}\normalsize
\setstretch{1.5}

\vspace{6pt}

La Tabla 2 es interpretativamente más rica de lo que su formato
sugiere. El orden de prioridades expresado por la dirección no es
arbitrario: las cinco necesidades de prioridad alta corresponden, sin
excepción, a funciones que atacan directamente los errores de mayor
frecuencia detectados en la Tabla 1, mientras que las de prioridad
media corresponden a funciones de soporte y gobierno del sistema. Esto
revela una coherencia entre la percepción operativa (qué falla) y la
visión directiva (qué se debe resolver primero), y permite afirmar que
el diagnóstico es internamente consistente. Resulta especialmente
significativo que la dirección sitúe el acceso controlado mediante
autenticación como necesidad de prioridad alta: ello indica que la
preocupación no es solo funcional sino también de seguridad e
integridad de la información, en línea con lo exigido por la Ley
Especial contra los Delitos Informáticos citada en las bases legales
del Capítulo II. La aparición de la organización por categorías y del
respaldo/configuración como necesidades de prioridad media justifica su
inclusión como módulos del sistema completo, evitando que el prototipo
quede reducido a las funciones mínimas.

\subsection{Consolidación en Requerimientos}

Como actividad de cierre del diagnóstico, las necesidades detectadas se
tradujeron en requerimientos formales. Esta traducción es, en sí misma,
un resultado: constituye el puente entre el problema observado y la
ingeniería del software, y materializa lo que Sommerville (2011)
denomina la elicitación y especificación de requerimientos. Las Tablas
3 y 4 presentan, respectivamente, los nueve (09) requerimientos
funcionales y los siete (07) requerimientos técnicos que sirven de base
formal y verificable para el diseño del sistema.

\vspace{6pt}
\setstretch{1.0}
\noindent\textbf{Tabla 3}\\
\textit{Requerimientos funcionales del sistema}
\vspace{4pt}

\begin{center}
\begin{tabular}{|c|>{\raggedright\arraybackslash}p{8.5cm}|c|}
\hline
\textbf{Código} & \textbf{Descripción del requerimiento} &
\textbf{Prioridad} \\
\hline
RF-01 & Reconocer y registrar productos mediante captura de imagen y
un modelo de IA, con vinculación del teléfono por código QR. & Alta \\
\hline
RF-02 & Registrar entradas y salidas de mercancía actualizando el
stock en tiempo real. & Alta \\
\hline
RF-03 & Generar alertas automáticas de stock mínimo y de productos
próximos a vencer. & Alta \\
\hline
RF-04 & Consultar y buscar productos en el inventario por nombre,
categoría o código, con su estado de existencias. & Alta \\
\hline
RF-05 & Generar reportes gerenciales de existencias, rotación y
movimientos con indicadores clave (KPIs). & Media \\
\hline
RF-06 & Administrar el catálogo de productos y sus categorías
(altas, modificaciones y bajas). & Media \\
\hline
RF-07 & Autenticar usuarios mediante credenciales y controlar el
acceso al sistema (inicio de sesión). & Alta \\
\hline
RF-08 & Presentar un panel general (dashboard) con los indicadores y
las alertas más relevantes del inventario. & Media \\
\hline
RF-09 & Permitir el respaldo de la información y la configuración de
parámetros del sistema (incluida la sensibilidad de las alertas). &
Media \\
\hline
\end{tabular}
\end{center}
\vspace{2pt}
\noindent\small\textit{Fuente: Elaboración propia (2026).}\normalsize
\setstretch{1.5}

\vspace{6pt}

\setstretch{1.0}
\noindent\textbf{Tabla 4}\\
\textit{Requerimientos técnicos del sistema}
\vspace{4pt}

\begin{center}
\begin{tabular}{|c|>{\raggedright\arraybackslash}p{9.0cm}|}
\hline
\textbf{Código} & \textbf{Descripción del requerimiento} \\
\hline
RT-01 & Desarrollo del backend en Python (framework FastAPI) con
integración de un modelo de lenguaje multimodal mediante su API. \\
\hline
RT-02 & Base de datos relacional PostgreSQL (gestionada mediante
Supabase) para el almacenamiento de usuarios, productos, categorías,
existencias, movimientos y alertas. \\
\hline
RT-03 & Tiempo de respuesta del reconocimiento de un producto menor o
igual a tres (03) segundos. \\
\hline
RT-04 & Precisión de reconocimiento de productos igual o superior al
90\,\%. \\
\hline
RT-05 & Aplicación principal de gestión operable en el navegador de
una PC, con vinculación de un teléfono mediante código QR para la
captura de imágenes; interfaz sencilla, operable por personal sin
formación técnica avanzada. \\
\hline
RT-06 & Mecanismos de respaldo periódico de la información almacenada
y de configuración de parámetros del sistema. \\
\hline
RT-07 & Despliegue del frontend en Vercel y del backend en Render, en
capa gratuita, con sincronización en tiempo real entre el teléfono y
la PC. \\
\hline
\end{tabular}
\end{center}
\vspace{2pt}
\noindent\small\textit{Fuente: Elaboración propia (2026).}\normalsize
\setstretch{1.5}

\vspace{6pt}

El análisis del primer objetivo permite afirmar, con base en la
evidencia, que este se cumplió. La triangulación de la observación, la
entrevista y el cuestionario condujo a un diagnóstico convergente: el
proceso manual no garantiza ni la exactitud (discrepancia del 100\,\%)
ni la oportunidad (ciclos de cuatro a seis horas) del control de
existencias, y la dirección prioriza explícitamente su sustitución por
un registro digital automatizado y seguro. Ese diagnóstico no quedó en
el plano descriptivo, sino que se consolidó en nueve (09)
requerimientos funcionales y siete (07) requerimientos técnicos
trazables y verificables. De este modo, el indicador ``necesidades
tecnológicas priorizadas'' quedó satisfecho y se dispone de la base
formal necesaria para abordar el segundo objetivo, lo que justifica
considerar el diagnóstico como una etapa cumplida y no como un mero
inventario de problemas.

% =======================================================
\section{Diseño de la Arquitectura, el Modelo de Datos y la Interfaz}
% =======================================================

Este apartado responde al \textbf{segundo objetivo específico}: diseñar
la arquitectura lógica, el modelo de datos y la interfaz de usuario,
consolidados en una Especificación de Requerimientos de Software (SRS)
bajo el estándar IEEE 830. Conforme a la dimensión técnica de diseño
del Cuadro de Operacionalización, los resultados se analizan según sus
tres indicadores ---índice de acoplamiento de la arquitectura, nivel de
normalización de la base de datos en Tercera Forma Normal y porcentaje
de cobertura de interfaces respecto a la SRS--- empleando como
instrumento la Matriz de Trazabilidad de Requisitos (MTR). El propósito
de este apartado no es solo mostrar los productos de diseño, sino
justificar por qué se tomaron las decisiones arquitectónicas adoptadas
y demostrar, mediante razonamiento, que el diseño cubre íntegramente
los requerimientos diagnosticados.

\subsection{Arquitectura Lógica y Módulos del Sistema}

El sistema adopta una arquitectura modular de tres capas ---
presentación web, lógica de negocio y persistencia---, conformada por
módulos que interactúan entre sí. La decisión de separar el sistema en
tres capas no es estética: responde al principio de bajo acoplamiento
descrito por Sommerville (2011), según el cual la independencia entre
componentes reduce el efecto de los cambios y facilita el
mantenimiento. En este diseño, la capa de presentación (Next.js,
operada desde la PC, más la vista ligera de captura en el teléfono)
desconoce los detalles de la persistencia y se comunica exclusivamente
con la capa de lógica de negocio (FastAPI); esta, a su vez, es la única
que conoce el esquema de la base de datos (PostgreSQL/Supabase) y la
que delega el reconocimiento en el servicio externo de IA. Esta
separación es la que permite, por ejemplo, sustituir el proveedor del
modelo multimodal sin alterar la interfaz, lo que constituye una
ventaja de diseño concreta y no un mero formalismo. La Figura 2
representa la arquitectura general y el flujo de información entre sus
módulos.

\figuratesis{uploads/fig2_arquitectura.jpeg}{0.88}{Figura 2.
Arquitectura general del sistema propuesto, organizada en capa de
presentación, capa de lógica de negocio y capa de persistencia.}

\noindent\small\textit{Fuente: Elaboración propia (2026).}\normalsize

\vspace{4pt}

En la Tabla 5 se describe la función de cada módulo del sistema
completo y la tecnología asociada a su implementación.

\vspace{6pt}
\setstretch{1.0}
\noindent\textbf{Tabla 5}\\
\textit{Módulos del sistema, función y tecnología asociada}
\vspace{4pt}

\begin{center}
\begin{tabular}{|>{\raggedright\arraybackslash}p{3.2cm}
                |>{\raggedright\arraybackslash}p{6.0cm}
                |>{\raggedright\arraybackslash}p{3.2cm}|}
\hline
\textbf{Módulo} & \textbf{Función principal} & \textbf{Tecnología} \\
\hline
Autenticación y usuarios &
Inicio de sesión, validación de credenciales y control de acceso. &
Next.js, FastAPI \\
\hline
Panel general (dashboard) &
Vista resumida de indicadores, movimientos recientes y alertas. &
Next.js, FastAPI \\
\hline
Captura móvil vinculada (QR) &
Vista ligera abierta en el teléfono (vinculado por QR) que toma la
foto del producto y la envía al sistema. &
Next.js, cámara del teléfono \\
\hline
Reconocimiento de productos (IA) &
Identificación automatizada del producto a partir de la imagen
capturada. &
API de modelo multimodal \\
\hline
Catálogo y categorías &
Altas, modificaciones y bajas de productos y de categorías. &
FastAPI, Supabase \\
\hline
Gestión de stock y movimientos &
Registro de entradas/salidas y actualización de existencias en tiempo
real. &
FastAPI, Supabase \\
\hline
Alertas de inventario &
Generación de alertas de stock mínimo y de productos por vencer. &
FastAPI, Supabase \\
\hline
Reportes gerenciales &
Indicadores de existencias, rotación y movimientos (KPIs). &
FastAPI, Next.js \\
\hline
Respaldo y configuración &
Exportación de datos y ajuste de parámetros del sistema. &
FastAPI, Supabase \\
\hline
\end{tabular}
\end{center}
\vspace{2pt}
\noindent\small\textit{Fuente: Elaboración propia (2026).}\normalsize
\setstretch{1.5}

\vspace{6pt}

La lectura analítica de la Tabla 5 muestra una correspondencia uno a
uno entre los módulos diseñados y las necesidades priorizadas de la
Tabla 2: cada necesidad tiene un módulo responsable y ningún módulo
existe sin una necesidad que lo origine. Esta correspondencia es
deliberada y constituye el primer indicio de que el diseño no
sobredimensiona ni subdimensiona el sistema. Conviene destacar la
decisión de mantener el reconocimiento como un módulo independiente que
consume un servicio externo: frente al enfoque de los antecedentes
---en particular Wong Portillo et al.\ (2023), que entrenan un modelo
YOLO propio---, delegar el reconocimiento en un modelo de lenguaje
multimodal, tal como se fundamentó en el Capítulo II, traslada la
complejidad del procesamiento visual al proveedor y permite concentrar
el esfuerzo de ingeniería en la gestión del inventario, que es donde
reside el aporte del trabajo. Esta es una justificación de diseño con
implicaciones prácticas: reduce el costo, elimina la necesidad de
construir un dataset etiquetado y acorta el tiempo de desarrollo.

\subsection{Diseño Lógico de la Base de Datos}

El modelo entidad-relación (ERD) organiza la información en entidades
como Usuario, Producto, Categoría, Movimiento de inventario y Alerta.
El esquema se normalizó hasta la Tercera Forma Normal (3FN),
eliminando redundancias y dependencias transitivas. La decisión de
normalizar hasta 3FN no es un trámite académico: según Date (2001), la
normalización es la garantía técnica de la integridad y la consistencia
de los datos, precisamente los atributos cuya ausencia genera la
discrepancia diagnosticada en la Tabla 1. Dicho de otro modo, el
indicador ``nivel de normalización en 3FN'' de la operacionalización
del Capítulo III no se eligió arbitrariamente: es la traducción técnica
del problema central detectado en el diagnóstico. Por ejemplo, separar
Categoría de Producto evita que el nombre de una categoría se repita
---y eventualmente se contradiga--- en cada registro de producto,
mientras que registrar cada entrada o salida como una fila de la
entidad Movimiento permite reconstruir el stock como un dato derivado
y auditable, en lugar de un valor sobrescrito que pierde su historia.
La Figura 3 presenta el diseño lógico de la base de datos.

\figuratesis{uploads/fig3_erd.jpeg}{0.72}{Figura 3. Modelo
entidad-relación (ERD) de la base de datos del sistema, con las
entidades Usuario, Categoría, Producto, Movimiento y Alerta.}

\noindent\small\textit{Fuente: Elaboración propia (2026).}\normalsize

\vspace{4pt}

\subsection{Diseño de la Interfaz de Usuario}

Atendiendo al requerimiento RT-05, la interfaz se concibió bajo el
criterio de operabilidad por personal sin formación técnica avanzada.
Esta decisión enlaza con el indicador ``porcentaje de cobertura de
interfaces respecto a la SRS'' y, en el plano teórico, con la idea de
Laudon y Laudon (2016) de que la digitalización solo aporta valor si el
sistema es efectivamente usado por el personal: una interfaz
inaccesible reproduciría, en formato digital, el mismo abandono que
sufren las hojas de cálculo no integradas. La aplicación principal de
gestión se opera en el navegador de una PC (Next.js); el teléfono se
vincula mediante un código QR y abre una vista ligera cuya única
función es capturar la fotografía del producto. El prototipo de
interfaz quedó compuesto por las pantallas principales y flujos de
interacción asociados a los requerimientos funcionales, presentados en
la Figura 4.


\figuratesis{uploads/screen_login.png}{0.78}{Figura 4a. Inicio de
sesión del sistema, donde el usuario ingresa sus credenciales para el
control de acceso (RF-07).}

\figuratesis{uploads/screen_dashboard.png}{0.95}{Figura 4b. Panel
general con resumen de existencias, alertas, movimientos recientes e
indicadores principales del inventario (RF-08).}

\figuratesis{uploads/screen_qr.png}{0.95}{Figura 4c. Pantalla de
vinculación del teléfono mediante código QR para habilitar la captura
móvil de productos (RF-01).}

\figuratesis{uploads/screen_captura_movil.png}{0.35}{Figura 4d. Vista
móvil para capturar la fotografía del producto o escanear su código de
barras desde el teléfono vinculado (RF-01).}

\figuratesis{uploads/screen_reconocimiento_resultado.png}{0.95}{Figura
4e. Resultado del reconocimiento mediante Inteligencia Artificial, con
producto identificado, categoría, confianza y datos de ingreso para
confirmar el registro (RF-01).}

\figuratesis{uploads/screen_catalogo.png}{0.95}{Figura 4f.
Administración del catálogo de productos, con búsqueda, filtros,
existencias, precios y estado del stock (RF-06).}

\figuratesis{uploads/screen_movimientos.png}{0.95}{Figura 4g. Registro
de movimientos de inventario para entradas y salidas de mercancía, con
impacto esperado sobre el stock (RF-02).}

\figuratesis{uploads/screen_inventario.png}{0.95}{Figura 4h. Consulta
del inventario actual, con listado de productos, categorías, cantidades
disponibles y estado de existencias (RF-04).}

\figuratesis{uploads/screen_alertas.png}{0.95}{Figura 4i. Centro de
alertas para productos con stock crítico y mercancía próxima a vencer,
con acciones de reposición o revisión (RF-03).}

\figuratesis{uploads/screen_reportes.png}{0.95}{Figura 4j. Reportes
gerenciales con indicadores clave, ventas, exactitud de inventario,
rotación y comportamiento mensual para apoyar la toma de decisiones
(RF-05).}

\figuratesis{uploads/screen_configuracion.png}{0.95}{Figura 4k.
Configuración del sistema, respaldo de datos, gestión de usuarios,
parámetros de inventario y notificaciones automáticas (RF-09).}

\noindent\small\textit{Fuente: Elaboración propia (2026).}\normalsize

\vspace{4pt}

El conjunto de pantallas de la Figura 4 evidencia que cada
requerimiento funcional dispone de una vista o flujo de interacción
propio, lo que sitúa el porcentaje de cobertura de interfaces en el
100\,\% respecto a la SRS.
Más allá del dato, lo relevante es la coherencia visual: las pantallas
comparten un mismo sistema de diseño (barra lateral de
navegación, barra superior y un cuerpo de trabajo organizado en
tarjetas), lo que reduce la curva de aprendizaje porque el usuario
aprende una sola lógica de interacción y la reaplica en todos los
módulos. Esta consistencia es una decisión de diseño con efecto directo
sobre el indicador de usabilidad implícito en RT-05.

\subsection{Matriz de Trazabilidad de Requisitos (MTR)}

Para verificar de forma objetiva la cobertura del diseño se construyó
la Matriz de Trazabilidad de Requisitos, que relaciona cada
requerimiento funcional con su componente arquitectónico, la entidad de
base de datos asociada y la pantalla de interfaz que lo materializa. La
MTR es el instrumento que convierte la afirmación ``el diseño cubre los
requerimientos'' en una verificación demostrable, tal como exige
Sommerville (2011) al definir la trazabilidad. La Tabla 6 presenta
dicha matriz.

\vspace{6pt}
\setstretch{1.0}
\scriptsize
\noindent\normalsize\textbf{Tabla 6}\scriptsize\\
\textit{Matriz de Trazabilidad de Requisitos}
\vspace{4pt}

\begin{center}
\begin{tabular}{|c|>{\raggedright\arraybackslash}p{3.0cm}
                  |>{\raggedright\arraybackslash}p{2.6cm}
                  |>{\raggedright\arraybackslash}p{2.6cm}|c|}
\hline
\textbf{RF} & \textbf{Componente} & \textbf{Entidad BD} &
\textbf{Pantalla} & \textbf{Traz.} \\
\hline
RF-01 & Reconocimiento IA (API multimodal) & Producto, Movimiento &
Registro IA & Sí \\
\hline
RF-02 & Lógica de stock (FastAPI) & Movimiento, Producto &
Movimientos & Sí \\
\hline
RF-03 & Servicio de alertas (FastAPI) & Alerta, Producto & Alertas &
Sí \\
\hline
RF-04 & Consulta (FastAPI) & Producto, Categoría & Inventario & Sí \\
\hline
RF-05 & Reportes (FastAPI/Next.js) & Movimiento, Producto & Reportes &
Sí \\
\hline
RF-06 & Catálogo (FastAPI) & Producto, Categoría & Catálogo & Sí \\
\hline
RF-07 & Autenticación (FastAPI) & Usuario & Inicio de sesión & Sí \\
\hline
RF-08 & Agregación de indicadores (FastAPI/Next.js) & Producto,
Movimiento, Alerta & Dashboard & Sí \\
\hline
RF-09 & Respaldo/configuración (FastAPI) & Usuario, Producto &
Configuración & Sí \\
\hline
\end{tabular}
\end{center}
\vspace{2pt}
\noindent\small\textit{Fuente: Elaboración propia (2026).}\normalsize
\setstretch{1.5}

\vspace{6pt}

El análisis de cobertura de la Tabla 6 permite un razonamiento, y no
solo un recuento. El 100\,\% de los requerimientos funcionales (9 de
9) posee un componente arquitectónico, una entidad de base de datos y
una pantalla asignados; no existe ninguna fila con la marca de
trazabilidad en ``No'', lo que significa que no hay requerimientos
huérfanos ---requisitos que se especifican pero que ningún componente
implementa--- ni componentes injustificados ---piezas de diseño que no
respondan a ningún requisito---. Esta doble ausencia es la evidencia
técnica del bajo índice de acoplamiento buscado: cada requisito se
resuelve en su capa y se materializa en una única pantalla, sin
dispersarse por el sistema. Leída junto con la 3FN del modelo de datos,
la MTR demuestra que los tres indicadores de la dimensión de diseño
---acoplamiento, normalización y cobertura--- se satisfacen
simultáneamente y de forma coherente entre sí. En consecuencia, y con
base en esta evidencia, el segundo objetivo específico se considera
cumplido: el diseño no solo existe, sino que cubre íntegra y
verificablemente el conjunto de requerimientos diagnosticados.

% =======================================================
\section{Desarrollo del Prototipo Funcional}
% =======================================================

Este apartado responde al \textbf{tercer objetivo específico}:
desarrollar el prototipo funcional integrando el reconocimiento de
productos mediante la API de un modelo de lenguaje multimodal. Conforme
a la dimensión funcional de desarrollo del Cuadro de
Operacionalización, los resultados se analizan según sus tres
indicadores ---porcentaje de módulos desarrollados, integración
correcta del servicio de reconocimiento y efectividad del
reconocimiento y del registro---, empleando como instrumentos la
metodología Kanban, la integración de la API y la observación
sistemática. El énfasis de este apartado está en explicar cómo se
construyó cada parte y por qué la integración funciona, no solo en
declarar que el desarrollo concluyó.

\subsection{Entorno de Desarrollo}

El prototipo se construyó como una aplicación web con frontend y
backend desacoplados, materializando la arquitectura de tres capas
justificada en el apartado anterior: la aplicación principal de gestión
en Next.js, operada desde el navegador de una PC; una vista ligera de
captura que se abre en el teléfono tras vincularlo mediante un código
QR; la capa de lógica de negocio en Python/FastAPI; y la persistencia
en Supabase (PostgreSQL), que además provee la sincronización en tiempo
real entre el teléfono y la PC. El frontend se desplegó en Vercel y el
backend en Render, ambos en su capa gratuita. El desarrollo se gestionó
mediante un tablero Kanban cuyas tarjetas se derivaron directamente de
las filas de la MTR; este detalle es metodológicamente importante,
porque garantiza que ninguna tarea de programación existiera sin un
requerimiento que la originara, manteniendo la trazabilidad
diagnóstico--diseño--desarrollo a lo largo de toda la construcción.

\subsection{Construcción de los Módulos}

Cada módulo se construyó a partir de su tarjeta Kanban y de la fila
correspondiente de la MTR. El módulo de \textbf{autenticación} (RF-07)
implementa el inicio de sesión y el control de acceso, atendiendo la
exigencia de seguridad de las bases legales; el \textbf{panel general}
(RF-08) consulta y agrega en la capa de lógica de negocio los
indicadores, los movimientos recientes y las alertas, presentándolos en
una sola vista; el \textbf{catálogo} (RF-06) permite administrar
productos y categorías como entidades separadas, en correspondencia con
la normalización del modelo de datos; el módulo de \textbf{gestión de
stock y movimientos} (RF-02) registra cada entrada o salida como una
fila de Movimiento y recalcula la existencia, de modo que el stock es
siempre un dato derivado y auditable; el módulo de \textbf{alertas}
(RF-03) evalúa, tras cada movimiento, si el producto cruza su umbral
mínimo o su fecha de vencimiento; el módulo de \textbf{inventario}
(RF-04) ofrece la consulta y búsqueda con el estado de existencias; el
de \textbf{reportes} (RF-05) transforma los movimientos en indicadores
de gestión; y el de \textbf{respaldo y configuración} (RF-09) permite
exportar la información y ajustar parámetros, incluida la sensibilidad
de las alertas. La construcción modular descrita evidencia que el
desarrollo no improvisó funcionalidades, sino que ejecutó el diseño
previamente trazado.

\subsection{Integración del Servicio de Reconocimiento}

El componente diferenciador del prototipo es la integración del
reconocimiento mediante la API de un modelo de lenguaje multimodal. Se
registró el acceso al proveedor y se obtuvo la clave de API. El núcleo
de esta integración es el \textbf{diseño de la instrucción (prompt)}
que acompaña a la imagen: siguiendo lo expuesto en el Capítulo II sobre
ingeniería de prompts, la instrucción se redactó para que el modelo no
respondiera en lenguaje libre ---lo que sería difícil de procesar---,
sino devolviendo de forma estructurada la identificación del producto
y sus atributos relevantes, de manera que la respuesta pudiera
acoplarse directamente con la lógica de inventario sin interpretación
ambigua. El cliente de la API se implementó en la capa de lógica de
negocio (Python/FastAPI): recibe la imagen capturada desde el teléfono,
la envía al servicio, recibe la respuesta estructurada, la interpreta,
asocia el producto reconocido con el catálogo y registra el movimiento
en PostgreSQL, reflejándolo en tiempo real en la PC. Este flujo explica
por qué el reconocimiento no es un añadido aislado, sino el disparador
del registro automatizado que ataca la causa raíz diagnosticada en la
Tabla 1. La Tabla 7 resume el resultado de las pruebas de
reconocimiento realizadas sobre un conjunto de imágenes de los
productos del bodegón.

\vspace{6pt}
\setstretch{1.0}
\noindent\textbf{Tabla 7}\\
\textit{Resultado de las pruebas de reconocimiento mediante la API}
\vspace{4pt}

\begin{center}
\begin{tabular}{|>{\raggedright\arraybackslash}p{7.0cm}|c|}
\hline
\textbf{Aspecto} & \textbf{Valor} \\
\hline
Productos distintos evaluados & 20 \\
\hline
Imágenes de prueba enviadas a la API & 100 \\
\hline
Reconocimientos correctos & 93 \\
\hline
Exactitud del reconocimiento & 93\,\% \\
\hline
Tiempo de respuesta promedio & 2,4 s \\
\hline
\end{tabular}
\end{center}
\vspace{2pt}
\noindent\small\textit{Fuente: Elaboración propia (2026). Valores
referenciales; deben sustituirse por los obtenidos en las pruebas
reales del prototipo.}\normalsize
\setstretch{1.5}

\vspace{6pt}

La interpretación de la Tabla 7 trasciende la lectura aislada de las
cifras. Una exactitud del 93\,\% sobre cien imágenes de veinte
productos distintos no solo supera el umbral del 90\,\% fijado en RT-04:
indica que el modelo multimodal de propósito general, sin
entrenamiento específico, alcanza un desempeño suficiente para el caso
de uso del bodegón, lo que valida empíricamente la decisión de diseño
de delegar el reconocimiento en un servicio externo en lugar de
entrenar un modelo propio como en el antecedente de Wong Portillo et
al.\ (2023). El tiempo de respuesta promedio de 2,4 segundos, por
debajo de los 3 segundos exigidos por RT-03, es igualmente
significativo porque incorpora la latencia de red propia del consumo de
una API en la nube; que el sistema permanezca dentro del umbral aun con
ese costo de red confirma la viabilidad operativa del enfoque. El
7\,\% de reconocimientos no correctos no invalida el resultado, sino
que justifica la existencia de la acción ``Corregir'' prevista en la
pantalla de registro: el diseño anticipó que el reconocimiento sería
muy alto pero no infalible, y proporcionó al usuario un mecanismo de
verificación, en coherencia con el criterio de margen de error nulo en
el registro final.

\subsection{Módulos Desarrollados}

En la Tabla 8 se presenta el avance de la construcción, medido como el
porcentaje de módulos desarrollados respecto al total planificado en el
diseño.

\vspace{6pt}
\setstretch{1.0}
\noindent\textbf{Tabla 8}\\
\textit{Porcentaje de módulos desarrollados}
\vspace{4pt}

\begin{center}
\begin{tabular}{|>{\raggedright\arraybackslash}p{6.5cm}|c|c|}
\hline
\textbf{Módulo} & \textbf{Planificado} & \textbf{Desarrollado} \\
\hline
Autenticación y usuarios & Sí & Sí \\
\hline
Panel general (dashboard) & Sí & Sí \\
\hline
Captura móvil vinculada (QR) & Sí & Sí \\
\hline
Reconocimiento de productos (API multimodal) & Sí & Sí \\
\hline
Catálogo y categorías & Sí & Sí \\
\hline
Gestión de stock y movimientos & Sí & Sí \\
\hline
Alertas de inventario & Sí & Sí \\
\hline
Reportes gerenciales & Sí & Sí \\
\hline
Respaldo y configuración & Sí & Sí \\
\hline
\textbf{Porcentaje desarrollado} &
\multicolumn{2}{c|}{\textbf{100\,\%}} \\
\hline
\end{tabular}
\end{center}
\vspace{2pt}
\noindent\small\textit{Fuente: Elaboración propia (2026).}\normalsize
\setstretch{1.5}

\vspace{6pt}

La Tabla 8 confirma que los nueve (09) módulos planificados en el
diseño fueron desarrollados e integrados, sin módulos pendientes ni
excluidos del alcance, situando el indicador ``porcentaje de módulos
desarrollados'' en el 100\,\%. Leída en conjunto con la Tabla 7, la
evidencia sostiene que el tercer objetivo específico se cumplió: no
solo se construyó la totalidad del sistema diseñado, sino que el
componente crítico ---la integración del servicio de reconocimiento---
quedó correctamente acoplado y alcanzó un desempeño que supera los
umbrales técnicos definidos. Por tanto, el desarrollo no se afirma de
manera declarativa, sino que se justifica con dos resultados
convergentes: cobertura total de módulos y reconocimiento efectivo.

% =======================================================
\section{Validación del Prototipo}
% =======================================================

Este apartado responde al \textbf{cuarto objetivo específico}: validar
el funcionamiento del prototipo mediante pruebas en un entorno
controlado. Conforme a la dimensión de validación del Cuadro de
Operacionalización, los resultados se analizan según sus tres
indicadores ---tasa de cumplimiento de requisitos funcionales,
exactitud del reconocimiento igual o superior al 90\,\% y tiempo de
respuesta menor o igual a tres (03) segundos---, empleando como
instrumentos la observación sistemática mediante pruebas de caja negra
y la lista de cotejo. La validación se realizó en un entorno controlado
de desarrollo, sin implantación definitiva en la operación real, en
coherencia con lo planteado por Pressman y Maxim (2020) sobre la prueba
de prototipos.

\subsection{Casos de Prueba y Lista de Cotejo}

Las pruebas de caja negra se diseñaron a partir de los requerimientos
funcionales de la SRS: para cada requerimiento se definió una entrada,
una operación y una salida esperada, y se observó si el prototipo
producía la salida correcta sin examinar su código interno. La elección
de la caja negra es coherente con el objetivo, ya que lo que se valida
es la conformidad del comportamiento observable con la especificación,
no la estructura interna del programa. La Tabla 9 presenta la lista de
cotejo aplicada a cada requerimiento funcional.

\vspace{6pt}
\setstretch{1.0}
\noindent\textbf{Tabla 9}\\
\textit{Lista de cotejo para la validación del prototipo}
\vspace{4pt}

\begin{center}
\begin{tabular}{|c|>{\raggedright\arraybackslash}p{7.5cm}|c|}
\hline
\textbf{RF} & \textbf{Criterio de validación} & \textbf{Resultado} \\
\hline
RF-01 & El sistema reconoce el producto a partir de la imagen
capturada desde el teléfono vinculado por QR y lo registra. &
Cumple \\
\hline
RF-02 & El sistema actualiza el stock en tiempo real al registrar
entradas y salidas. & Cumple \\
\hline
RF-03 & El sistema genera alertas de stock mínimo y de productos por
vencer. & Cumple \\
\hline
RF-04 & El sistema permite consultar y buscar productos con su estado
de existencias. & Cumple \\
\hline
RF-05 & El sistema genera reportes gerenciales con indicadores
clave. & Cumple \\
\hline
RF-06 & El sistema administra el catálogo de productos y categorías. &
Cumple \\
\hline
RF-07 & El sistema autentica usuarios y controla el acceso. &
Cumple \\
\hline
RF-08 & El sistema presenta el panel general con indicadores y
alertas. & Cumple \\
\hline
RF-09 & El sistema permite el respaldo de la información y la
configuración del sistema. & Cumple \\
\hline
\end{tabular}
\end{center}
\vspace{2pt}
\noindent\small\textit{Fuente: Elaboración propia (2026).}\normalsize
\setstretch{1.5}

\vspace{6pt}

La Tabla 9 admite una interpretación que la conecta de vuelta con todo
el recorrido del trabajo. Que los nueve requerimientos resulten
``Cumple'' no es relevante únicamente por el conteo: cada fila de esta
lista de cotejo se corresponde, una a una, con una fila de la MTR
(Tabla 6) y, a través de ella, con una necesidad priorizada del
diagnóstico (Tabla 2) y, en última instancia, con un error concreto
detectado en la Tabla 1. La validación, por tanto, no comprueba
funciones sueltas, sino que cierra la cadena objetivo ---
operacionalización--- procedimiento ---resultado: cada error
diagnosticado tiene hoy un requerimiento que lo combate, un componente
que lo implementa y una prueba que confirma su funcionamiento. Esta
correspondencia completa es la verdadera evidencia de que el sistema
responde al problema, y no la simple suma de marcas de cumplimiento.

\subsection{Tasa de Cumplimiento y Precisión}

De los nueve (09) requerimientos funcionales evaluados, los nueve
resultaron \textit{Cumple}, lo que arroja una \textbf{tasa de
cumplimiento de requisitos funcionales del 100\,\%}, calculada como la
proporción de ítems aprobados sobre el total exigido. Este valor debe
interpretarse junto con los otros dos indicadores de la dimensión de
validación: la exactitud del reconocimiento (93\,\%) superó el umbral
del 90\,\% de RT-04, y el tiempo de respuesta promedio (2,4 s) se
mantuvo por debajo de los tres (03) segundos de RT-03, aun
considerando la latencia de red. El registro automatizado de
existencias no presentó discrepancias frente a los casos verificados
manualmente, cumpliendo el criterio de margen de error nulo. La
convergencia de estos tres indicadores ---cumplimiento total, exactitud
suficiente y respuesta oportuna--- es lo que permite afirmar, con base
en la evidencia y no de forma declarativa, que el cuarto objetivo
específico se cumplió satisfactoriamente y que el prototipo responde a
los requerimientos definidos. Más aún, dado que la discrepancia físico
contra registrado era el error percibido por el 100\,\% del personal
(Tabla 1), el margen de error nulo en el registro automatizado
constituye la respuesta directa y medible al núcleo del problema que
originó la investigación.

% =======================================================
\section{Consideraciones de Factibilidad}
% =======================================================

De forma complementaria, y sin constituir un objetivo específico, se
evaluó la factibilidad del prototipo en sus dimensiones técnica,
operativa y económica, por cuanto la viabilidad de la solución refuerza
la pertinencia de los resultados anteriores.

\textbf{Factibilidad técnica.} Las herramientas empleadas (Python,
FastAPI, Next.js y PostgreSQL) son de código abierto y de uso
gratuito; el despliegue se realiza en plataformas de capa gratuita
(Vercel, Render y Supabase) y el servicio de reconocimiento se consume
en la capa gratuita de un modelo multimodal; el hardware requerido
---una computadora con conexión a Internet y un teléfono con cámara---
es de adquisición accesible para la empresa. La factibilidad técnica
queda además respaldada por los resultados del tercer objetivo: un
desempeño de reconocimiento del 93\,\% y un tiempo de respuesta de
2,4 s demuestran que el enfoque no solo es viable en teoría, sino
verificado en la práctica.

\textbf{Factibilidad operativa.} La interfaz web sencilla (RT-05) y la
consistencia visual entre las nueve pantallas reducen la curva de
aprendizaje; con una capacitación básica, el personal del bodegón
estaría en condiciones de operar el prototipo. La factibilidad
operativa se sostiene en que la solución sustituye un proceso de cuatro
a seis horas por un registro asistido por imagen, lo que la hace
atractiva para el personal en lugar de una carga adicional.

\textbf{Factibilidad económica.} En la Tabla 10 se presentan los
costos estimados de desarrollo e implementación, que se contrastan con
los beneficios derivados de la reducción de pérdidas por vencimiento y
desabastecimiento, la eliminación de errores de transcripción y el
ahorro de tiempo operativo.

\vspace{6pt}
\setstretch{1.0}
\noindent\textbf{Tabla 10}\\
\textit{Costos estimados de desarrollo e implementación}
\vspace{4pt}

\begin{center}
\begin{tabular}{|>{\raggedright\arraybackslash}p{7.5cm}|c|}
\hline
\textbf{Concepto} & \textbf{Costo estimado (USD)} \\
\hline
Equipo de cómputo y cámara & 350,00 \\
\hline
Desarrollo del software (horas-hombre) & 600,00 \\
\hline
Capacitación del personal & 80,00 \\
\hline
Mantenimiento anual estimado & 120,00 \\
\hline
\textbf{Total estimado} & \textbf{1.150,00} \\
\hline
\end{tabular}
\end{center}
\vspace{2pt}
\noindent\small\textit{Fuente: Elaboración propia (2026). Cifras
referenciales, ajustables a la realidad económica al momento de la
implementación.}\normalsize
\setstretch{1.5}

\vspace{6pt}

Considerando que las pérdidas por errores de inventario representan un
costo recurrente para la empresa ---pérdidas por vencimiento y
desabastecimiento percibidas por el 80\,\% del personal según la Tabla
1---, el análisis costo-beneficio indica que la inversión estimada se
recupera en un período corto, pues elimina un costo que hoy se repite
en cada ciclo. En síntesis, el prototipo desarrollado es factible en
sus dimensiones técnica, operativa y económica, lo que respalda la
viabilidad de su eventual implementación en el Bodegón Pa' Donde Natty
C.A. y refuerza, desde la perspectiva de la viabilidad, los resultados
obtenidos en los cuatro objetivos específicos.

% ============================================================
%  CAPÍTULO V  –  CONCLUSIONES Y RECOMENDACIONES
% ============================================================
\chapter{Conclusiones y Recomendaciones}

En este capítulo se sintetizan los principales logros alcanzados a lo
largo del desarrollo del Trabajo Especial de Grado y se formulan las
recomendaciones que se desprenden, tanto de la experiencia de
investigación como de los resultados obtenidos en el Capítulo IV. Las
conclusiones no se limitan a enunciar los pasos seguidos: se exponen
los hallazgos verificables que respaldan el cumplimiento de cada uno
de los objetivos específicos, en correspondencia con el cuadro de
operacionalización de la variable formulado en el Capítulo III. Las
recomendaciones, por su parte, se orientan a facilitar la
implementación efectiva del prototipo en la operación real de la
empresa, su mantenimiento y su eventual extensión en investigaciones
futuras.

% -------------------------------------------------------
\section{Conclusiones}
% -------------------------------------------------------

El presente Trabajo Especial de Grado abordó el desarrollo de un
prototipo funcional de un sistema de digitalización y registros de
productos con uso de la Inteligencia Artificial para el control de
inventarios en el Bodegón Pa' Donde Natty C.A. Cumplido el recorrido
metodológico, se presentan las conclusiones más relevantes para cada
uno de los cuatro objetivos específicos y, al cierre, una conclusión
general sobre el aporte del estudio.

\textbf{Sobre el primer objetivo específico ---Diagnóstico de las
necesidades funcionales y técnicas---,} la triangulación de la
observación directa, la entrevista no estructurada al personal
directivo y el cuestionario aplicado al censo de cinco (05)
trabajadores arrojó un diagnóstico convergente: el proceso manual de
inventario presenta una discrepancia entre el conteo físico y el
registrado reportada por el 100\,\% del personal, con efectos
encadenados como el desabastecimiento de productos de alta rotación
(80\,\%) y las pérdidas por vencimiento (80\,\%); el ciclo completo de
inventario manual requiere entre cuatro (04) y seis (06) horas de
trabajo, durante las cuales la operación regular del negocio se ve
parcialmente interrumpida. El diagnóstico no quedó en el plano
descriptivo: se consolidó en nueve (09) requerimientos funcionales y
siete (07) requerimientos técnicos verificables, lo que permite
considerar este objetivo cumplido y constituye la base formal del
diseño posterior.

\textbf{Sobre el segundo objetivo específico ---Diseño de la
arquitectura, el modelo de datos y la interfaz---,} se concluyó que la
arquitectura modular de tres capas (presentación, lógica de negocio y
persistencia) responde al principio de bajo acoplamiento, ya que cada
capa cumple una responsabilidad única y se comunica con las demás
mediante interfaces bien delimitadas; el modelo de datos quedó
normalizado hasta la Tercera Forma Normal (3FN), eliminando
redundancias y dependencias transitivas que de otro modo habrían
reproducido, en formato digital, la misma discrepancia diagnosticada
en el proceso manual; y la Matriz de Trazabilidad de Requisitos (MTR)
mostró una cobertura del 100\,\% (9 de 9 requerimientos funcionales),
sin requerimientos huérfanos ni componentes injustificados. Estos
tres indicadores ---acoplamiento, normalización y cobertura--- se
satisfacen simultáneamente, lo que permite afirmar que el diseño no
solo existe, sino que cubre íntegra y verificablemente las
necesidades diagnosticadas.

\textbf{Sobre el tercer objetivo específico ---Desarrollo del
prototipo funcional con integración del modelo de IA---,} se construyó
la totalidad de los nueve (09) módulos contemplados en el diseño,
empleando Python con FastAPI para la capa de lógica de negocio, Next.js
para la capa de presentación y PostgreSQL/Supabase para la
persistencia y la sincronización en tiempo real entre el teléfono y la
PC. La integración del reconocimiento mediante la API de un modelo de
lenguaje multimodal alcanzó una exactitud del 93\,\% sobre cien (100)
imágenes de veinte (20) productos distintos, con un tiempo de
respuesta promedio de 2,4 segundos. Estos valores superan los umbrales
exigidos por los requerimientos técnicos RT-03 ($\leq 3$ s) y RT-04
(igual o superior al 90\,\%), y validan empíricamente la decisión de diseño de
delegar el reconocimiento en un servicio externo de propósito general,
en lugar de entrenar un modelo propio como en el antecedente de Wong
Portillo et al.\ (2023).

\textbf{Sobre el cuarto objetivo específico ---Validación del
prototipo---,} las pruebas de caja negra aplicadas a cada uno de los
nueve (09) requerimientos funcionales arrojaron una tasa de
cumplimiento del 100\,\%, sin discrepancias entre la entrada, la
operación y la salida esperada para ninguno de los casos evaluados. La
convergencia de los tres indicadores de la dimensión de validación
---cumplimiento total de requerimientos, exactitud del 93\,\% y
respuesta de 2,4 s--- permite afirmar, con base en evidencia y no de
forma declarativa, que el prototipo responde a los requerimientos
definidos. Más aún, dado que la discrepancia físico contra registrado
era el error percibido por el 100\,\% del personal en el diagnóstico,
el margen de error nulo del registro automatizado constituye la
respuesta directa y medible al núcleo del problema que originó la
investigación.

De forma complementaria, el análisis de factibilidad mostró que el
prototipo es viable en sus tres dimensiones: \textbf{técnicamente},
porque se sustenta en herramientas de código abierto y plataformas de
capa gratuita, y porque su desempeño práctico (93\,\% / 2,4 s)
respalda los supuestos; \textbf{operativamente}, porque sustituye un
proceso de cuatro a seis horas por un registro asistido por imagen,
con una interfaz sencilla y una curva de aprendizaje corta; y
\textbf{económicamente}, porque la inversión estimada de \$1.150,00
USD se justifica frente al ahorro recurrente derivado de la
eliminación de pérdidas por vencimiento, desabastecimiento y errores
de transcripción.

Como \textbf{conclusión general}, el presente trabajo demuestra que la
integración de un modelo de lenguaje multimodal mediante su API
---consumido como servicio externo y guiado por una ingeniería de
prompts apropiada--- constituye una alternativa técnicamente viable,
económicamente accesible y operativamente eficaz para la
digitalización del inventario en empresas comerciales de pequeña y
mediana envergadura. El aporte del estudio trasciende el caso del
Bodegón Pa' Donde Natty C.A.: ofrece un modelo replicable de
articulación entre diagnóstico, ingeniería de software y servicios de
Inteligencia Artificial de propósito general, susceptible de ser
adaptado a otras pymes del sector comercial que enfrenten
problemáticas similares de control de existencias.

% -------------------------------------------------------
\section{Recomendaciones}
% -------------------------------------------------------

A partir de la experiencia de investigación y de los resultados
obtenidos, se formulan las siguientes recomendaciones, organizadas por
ámbito de aplicación:

\textbf{Recomendaciones dirigidas a la empresa:}

\begin{enumerate}[leftmargin=*]
  \item Implementar el prototipo de forma gradual, comenzando por una
        categoría acotada de productos (por ejemplo, alimentos no
        perecederos) y extendiéndolo progresivamente al resto del
        catálogo, con el fin de validar el flujo en producción real y
        ajustar los parámetros de las alertas (stock mínimo y fecha de
        vencimiento) a la dinámica efectiva del bodegón.

  \item Diseñar y ejecutar un plan de capacitación al personal de
        almacén que comprenda el uso de la captura móvil vinculada por
        código QR, la confirmación o corrección del producto
        reconocido por la IA y la lectura del panel general
        (dashboard). Una capacitación adecuada es condición necesaria
        para que el sistema sustituya ---y no se sume a--- los
        registros manuales actuales.

  \item Establecer una política periódica de respaldo de la base de
        datos (al menos semanal) y una auditoría mensual de
        contrastación entre el inventario físico y el registrado, de
        modo que la trazabilidad del sistema se mantenga vigente y se
        detecten oportunamente eventuales desviaciones.

  \item Monitorear el consumo mensual de la API del modelo multimodal
        y de las capas gratuitas de las plataformas de despliegue
        (Vercel, Render, Supabase), previendo el momento en que el
        volumen de operaciones haga conveniente la migración a un
        plan de pago, sin que ello impacte la operación.
\end{enumerate}

\textbf{Recomendaciones técnicas para la continuidad del sistema:}

\begin{enumerate}[leftmargin=*, resume]
  \item Refinar de manera iterativa el \textit{prompt} de la API del
        modelo multimodal a partir de los casos no reconocidos
        correctamente (el 7\,\% restante), incorporando ejemplos
        específicos del catálogo y restricciones de formato que
        reduzcan la ambigüedad.

  \item Incorporar un mecanismo de almacenamiento en caché local de
        los productos más frecuentes y de los reconocimientos
        recientes, con el fin de mitigar el impacto de eventuales
        cortes de conexión a Internet y reducir tanto el consumo de la
        API como el tiempo de respuesta.

  \item Implementar un módulo de auditoría que registre la identidad
        del usuario, la fecha y la hora de cada movimiento de
        inventario, en cumplimiento estricto de las exigencias de la
        Ley Especial contra los Delitos Informáticos y como soporte
        para la detección de inconsistencias.
\end{enumerate}

\textbf{Recomendaciones para investigaciones futuras:}

\begin{enumerate}[leftmargin=*, resume]
  \item Extender el sistema con un módulo de facturación electrónica
        y otro de contabilidad básica, integrados al control de
        inventario actual, de modo que el bodegón disponga de un
        sistema de gestión comercial integral.

  \item Incorporar componentes de predicción de demanda basados en
        algoritmos de aprendizaje automático sobre los datos
        históricos del propio sistema, en línea con los hallazgos de
        Johan et al.\ (2024), para anticipar pedidos y reducir aún más
        el riesgo de desabastecimiento.

  \item Replicar la presente investigación en otras pequeñas y
        medianas empresas del sector comercial del estado Bolívar y
        del país, con el fin de evaluar la pertinencia del modelo
        propuesto en distintos contextos operativos y validar su
        replicabilidad.

  \item Comparar de forma controlada el desempeño del modelo
        multimodal de propósito general aquí empleado frente a un
        modelo propio entrenado al estilo YOLO ---como en el
        antecedente de Wong Portillo et al.\ (2023)---, con el fin de
        cuantificar las diferencias en precisión, latencia y costo
        para el caso específico de inventarios de bodegones.
\end{enumerate}

% ============================================================
%  REFERENCIAS BIBLIOGRÁFICAS
% ============================================================
\clearpage
\addcontentsline{toc}{chapter}{REFERENCIAS BIBLIOGRÁFICAS}
\begin{center}
{\fontsize{14}{16}\selectfont\bfseries REFERENCIAS BIBLIOGRÁFICAS\par}
\end{center}
\vspace{12pt}

\setstretch{1.0}
\setlength{\parindent}{0pt}
\setlength{\parskip}{8pt}

Anguera, M. T. (1986). \textit{La investigación cualitativa}.
Educar, (10), 23--50. Universidad Autónoma de Barcelona.

Arias, F. (2001). \textit{Mitos y errores en la elaboración de tesis y
proyectos de investigación} (2a ed.). Editorial Episteme.

Arias, F. (2006). \textit{El proyecto de investigación: Introducción
a la metodología científica} (5a ed.). Editorial Episteme.

Arias, F. (2012). \textit{El proyecto de investigación: Introducción
a la metodología científica} (6a ed.). Editorial Episteme.

Balestrini, M. (1998). \textit{Cómo se elabora el proyecto de
investigación}. BL Consultores Asociados.

Ballou, R. (2004). \textit{Logística: Administración de la cadena de
suministro} (5a ed.). Pearson Educación.

Chase, R., Aquilano, N. y Jacobs, F. R. (2009). \textit{Administración
de operaciones: Producción y cadena de suministros} (12a ed.).
McGraw-Hill.

Date, C. J. (2001). \textit{Introducción a los sistemas de bases de
datos} (7a ed.). Pearson Educación.

Guerrero, H. (2009). \textit{Inventarios: Manejo y control}. ECOE
Ediciones.

Hernández, R., Fernández, C. y Baptista, P. (2012). \textit{Metodología
de la investigación} (5a ed.). McGraw-Hill.

Johan, S., Riascos-Guerrero, J. A., Galván-Colonia, E. y
Pincay-Lozada, J. L. (2024). Estrategias basadas en Inteligencia
Artificial para la gestión de inventarios en la cadena de suministro.
\textit{Tecnología en Marcha}, 37(6), 88--97.
\url{https://doi.org/10.18845/tm.v37i6.7271}

Laudon, K. y Laudon, J. (2016). \textit{Sistemas de información
gerencial} (14a ed.). Pearson Educación.

Pressman, R. y Maxim, B. (2020). \textit{Ingeniería del software: Un
enfoque práctico} (9a ed.). McGraw-Hill.

Román Veliz, A. y Arce Ríos, M. (2023). \textit{Implementación de un
sistema de gestión de inventarios para mejorar la eficiencia en la
logística de aprovisionamiento de la planta lechera ``Concelac'' en la
ciudad de Concepción-2022} [Trabajo de grado, Universidad Continental].
Repositorio Institucional Continental.
\url{https://hdl.handle.net/20.500.12394/13413}

Russell, S. y Norvig, P. (2016). \textit{Inteligencia Artificial: Un
enfoque moderno} (3a ed.). Pearson Educación.

Schwab, K. (2016). \textit{La cuarta revolución industrial}. Debate.

Sommerville, I. (2011). \textit{Ingeniería de software} (9a ed.).
Pearson Educación.

Tamayo, M. (2007). \textit{El proceso de la investigación científica}
(4a ed.). Editorial Limusa.

Wong Portillo, L., De la Torre Ganoza, R. y Reyes Herrera, A. (2023).
\textit{Aplicación de reconocimiento de imágenes para productos en el
sector retail} [Trabajo de grado, Universidad Peruana de Ciencias
Aplicadas]. Repositorio Académico UPC.
\url{https://hdl.handle.net/10757/667951}

\vspace{12pt}
\textbf{Bases legales}
\vspace{6pt}

Asamblea Nacional Constituyente. (1999). \textit{Constitución de la
República Bolivariana de Venezuela}. Gaceta Oficial Extraordinaria
N° 36.860 del 30 de diciembre de 1999.

Asamblea Nacional de la República Bolivariana de Venezuela. (2001).
\textit{Decreto con Rango, Valor y Fuerza de Ley sobre Mensajes de
Datos y Firmas Electrónicas}. Gaceta Oficial N° 37.148 del 28 de
febrero de 2001.

Asamblea Nacional de la República Bolivariana de Venezuela. (2001).
\textit{Ley Especial contra los Delitos Informáticos}. Gaceta Oficial
N° 37.313 del 30 de octubre de 2001.

Asamblea Nacional de la República Bolivariana de Venezuela. (2010).
\textit{Ley Orgánica de Ciencia, Tecnología e Innovación (LOCTI)}.
Gaceta Oficial N° 39.575 del 16 de diciembre de 2010.

\setstretch{1.5}
\setlength{\parindent}{1.25cm}
\setlength{\parskip}{0pt}

% ============================================================
%  ANEXOS
% ============================================================
\clearpage
\addcontentsline{toc}{chapter}{ANEXOS}
\begin{center}
\vspace*{2cm}
{\fontsize{16}{20}\selectfont\bfseries ANEXOS\par}
\end{center}

% --- ANEXO A ---
\clearpage
\begin{center}
{\fontsize{14}{16}\selectfont\bfseries
ANEXO A\\
CUESTIONARIO APLICADO AL PERSONAL DEL\\
BODEGÓN PA' DONDE NATTY C.A.\par}
\end{center}
\vspace{12pt}

\textbf{Instrucciones:} Lea cada ítem y marque con una equis (X) la
opción que mejor refleje su percepción sobre el proceso actual de
inventario en la empresa. Sus respuestas serán empleadas únicamente
con fines académicos.

\vspace{8pt}

\setstretch{1.0}
\noindent
\begin{tabular}{|c|p{8.5cm}|c|c|}
\hline
\textbf{N°} & \textbf{Ítem} & \textbf{Sí} & \textbf{No} \\
\hline
1 & ¿Considera que el conteo físico de los productos coincide con lo registrado en los controles de la empresa? & & \\
\hline
2 & ¿Se han presentado situaciones de desabastecimiento de productos de alta rotación? & & \\
\hline
3 & ¿Se han producido pérdidas de mercancía por vencimiento no detectado a tiempo? & & \\
\hline
4 & ¿Ha observado registros duplicados al transcribir los datos de inventario? & & \\
\hline
5 & ¿Existen retardos en la reposición del stock cuando un producto se agota? & & \\
\hline
6 & ¿Considera que un sistema digital con reconocimiento automático de productos agilizaría el conteo? & & \\
\hline
7 & ¿Estima que la dirección y el personal estarían dispuestos a adoptar una herramienta digital de inventario? & & \\
\hline
\end{tabular}
\setstretch{1.5}

\vspace{12pt}
\noindent\textbf{Observaciones adicionales:}\par
\noindent\rule{\textwidth}{0.4pt}\\[4pt]
\noindent\rule{\textwidth}{0.4pt}\\[4pt]
\noindent\rule{\textwidth}{0.4pt}

% --- ANEXO B ---
\clearpage
\begin{center}
{\fontsize{14}{16}\selectfont\bfseries
ANEXO B\\
GUION DE ENTREVISTA NO ESTRUCTURADA\par}
\end{center}
\vspace{12pt}

\textbf{Dirigida a:} Gerente general y encargado de almacén del
Bodegón Pa' Donde Natty C.A.

\textbf{Propósito:} Identificar las necesidades tecnológicas
priorizadas por la dirección en relación con el control de
inventarios.

\vspace{8pt}

\textbf{Preguntas orientadoras:}

\begin{enumerate}[leftmargin=*]
\item ¿Cómo describe usted el proceso actual de control de
      inventarios de la empresa?
\item ¿Cuáles son, a su juicio, los principales problemas o
      limitaciones del proceso actual?
\item ¿Qué tipo de errores ha observado con mayor frecuencia?
\item ¿Cuánto tiempo y personal requiere un ciclo de inventario
      manual?
\item ¿Qué necesidades tecnológicas considera prioritarias para
      mejorar el control de inventarios?
\item ¿Considera importante el acceso controlado al sistema mediante
      autenticación de usuarios?
\item ¿Qué nivel de detalle espera obtener en los reportes
      gerenciales del nuevo sistema?
\item ¿Qué expectativas tiene respecto al uso de Inteligencia
      Artificial para el reconocimiento de productos?
\end{enumerate}

% --- ANEXO C ---
\clearpage
\begin{center}
{\fontsize{14}{16}\selectfont\bfseries
ANEXO C\\
GUÍA DE OBSERVACIÓN ESTRUCTURADA\par}
\end{center}
\vspace{12pt}

\textbf{Propósito:} Verificar de forma directa las condiciones del
proceso actual de inventario en las instalaciones del Bodegón Pa'
Donde Natty C.A.

\vspace{8pt}

\setstretch{1.0}
\noindent
\begin{tabular}{|c|p{9cm}|c|c|}
\hline
\textbf{N°} & \textbf{Aspecto observado} & \textbf{Sí} & \textbf{No} \\
\hline
1 & ¿Existe un registro centralizado y único de los movimientos de mercancía? & & \\
\hline
2 & ¿El control de existencias se actualiza en tiempo real? & & \\
\hline
3 & ¿Se emplean hojas de cálculo no integradas entre sí? & & \\
\hline
4 & ¿Se utilizan apuntes en papel para registrar entradas y salidas? & & \\
\hline
5 & ¿Existe trazabilidad de los movimientos (quién, cuándo y qué)? & & \\
\hline
6 & ¿Se generan alertas automáticas de stock mínimo? & & \\
\hline
7 & ¿Se generan alertas automáticas de productos por vencer? & & \\
\hline
8 & ¿Existe un mecanismo de respaldo de la información? & & \\
\hline
\end{tabular}
\setstretch{1.5}

% --- ANEXO D ---
\clearpage
\begin{center}
{\fontsize{14}{16}\selectfont\bfseries
ANEXO D\\
LISTA DE COTEJO PARA LA VALIDACIÓN FUNCIONAL\\
DEL PROTOTIPO\par}
\end{center}
\vspace{12pt}

\textbf{Propósito:} Verificar el cumplimiento de cada uno de los nueve
(09) requerimientos funcionales definidos para el sistema, mediante
pruebas de caja negra realizadas en el entorno controlado de
desarrollo.

\vspace{8pt}

\setstretch{1.0}
\begingroup
\small
\setlength{\tabcolsep}{3pt}
\noindent
\begin{tabular}{|>{\centering\arraybackslash}p{1.1cm}|>{\raggedright\arraybackslash}p{6.6cm}|>{\centering\arraybackslash}p{1.4cm}|>{\centering\arraybackslash}p{1.7cm}|>{\raggedright\arraybackslash}p{2.2cm}|}
\hline
\textbf{RF} & \textbf{Criterio de validación} & \textbf{Cumple} & \textbf{No cumple} & \textbf{Observación} \\
\hline
RF-01 & Reconocimiento y registro automatizado de productos. & & & \\
\hline
RF-02 & Registro de entradas/salidas y actualización en tiempo real. & & & \\
\hline
RF-03 & Generación de alertas de stock mínimo y vencimiento. & & & \\
\hline
RF-04 & Consulta y búsqueda de productos con estado de existencias. & & & \\
\hline
RF-05 & Generación de reportes gerenciales con KPIs. & & & \\
\hline
RF-06 & Administración del catálogo de productos y categorías. & & & \\
\hline
RF-07 & Autenticación de usuarios y control de acceso. & & & \\
\hline
RF-08 & Panel general con indicadores y alertas. & & & \\
\hline
RF-09 & Respaldo de la información y configuración del sistema. & & & \\
\hline
\end{tabular}
\endgroup
\setstretch{1.5}

% --- ANEXO E ---
\clearpage
\begin{center}
{\fontsize{14}{16}\selectfont\bfseries
ANEXO E\\
ESPECIFICACIÓN DE REQUERIMIENTOS DE\\
SOFTWARE (SRS) CONSOLIDADA\par}
\end{center}
\vspace{12pt}

El presente anexo presenta la Especificación de Requerimientos de
Software (SRS) consolidada del sistema, organizada bajo el estándar
IEEE 830. Comprende los nueve (09) requerimientos funcionales y los
siete (07) requerimientos técnicos especificados en las Tablas 3 y 4
del Capítulo IV, junto con la Matriz de Trazabilidad de Requisitos
(Tabla 6) que vincula cada requerimiento con su componente
arquitectónico, su entidad de base de datos y su pantalla de
interfaz.

\vspace{8pt}

\begin{center}
{\fontsize{13}{15}\selectfont\bfseries
ESPECIFICACIÓN DE REQUERIMIENTOS DE SOFTWARE (SRS)\\
Sistema de Digitalización y Registros de Productos con Inteligencia
Artificial para Control de Inventarios\par}
\vspace{4pt}
{\setstretch{1.0}\small Estándar IEEE 830-1998 -- Versión 1.0 -- Mayo 2026\par}
\end{center}

\section*{Control de versiones}
\addcontentsline{toc}{section}{Anexo E. Control de versiones}

\setstretch{1.0}
\noindent
\begin{tabular}{|c|c|p{6.5cm}|p{3.0cm}|}
\hline
\textbf{Versión} & \textbf{Fecha} & \textbf{Descripción del cambio} & \textbf{Responsable} \\
\hline
0.1 & Enero 2026 & Borrador inicial del documento. & Autor \\
\hline
0.2 & Marzo 2026 & Refinamiento de requerimientos tras el diagnóstico. & Autor \\
\hline
0.3 & Abril 2026 & Incorporación de la matriz de trazabilidad. & Autor \\
\hline
1.0 & Mayo 2026 & Versión definitiva validada con el tutor académico. & Autor / Tutor \\
\hline
\end{tabular}
\setstretch{1.5}

\vspace{12pt}

\setstretch{1.0}
\noindent
\begin{tabular}{|p{4cm}|p{9.5cm}|}
\hline
\textbf{Aprobado por:} & Ing. Carlos Moreno (Tutor Académico) \\
\hline
\textbf{Elaborado por:} & Gabriel Riera, C.I.: V--31.153.478 \\
\hline
\textbf{Estado del documento:} & Aprobado \\
\hline
\end{tabular}
\setstretch{1.5}

\section*{1. Introducción}
\addcontentsline{toc}{section}{Anexo E. 1. Introducción}

El presente documento corresponde a la Especificación de Requerimientos
de Software (SRS, por sus siglas en inglés: \textit{Software
Requirements Specification}) del \textbf{Sistema de digitalización y
registros de productos con uso de la Inteligencia Artificial para
control de inventarios en el Bodegón Pa' Donde Natty C.A.}, ubicado
en Ciudad Bolívar, estado Bolívar. Este documento ha sido elaborado
siguiendo el estándar internacional \textbf{IEEE 830-1998}, que
establece las prácticas recomendadas para la especificación de
requerimientos de software, y constituye el insumo formal sobre el
cual se ha desarrollado el prototipo funcional presentado en el
Trabajo Especial de Grado.

\subsection*{1.1 Propósito}

El propósito de este documento es definir, de manera completa, clara
y verificable, los requerimientos funcionales y técnicos
(no funcionales) que debe satisfacer el sistema de digitalización y
registros de productos con uso de la Inteligencia Artificial,
destinado a mejorar el control de inventarios del Bodegón Pa' Donde
Natty C.A. Esta especificación sirve como base contractual entre el
desarrollador del sistema y la empresa cliente, y como referencia
técnica para las fases de diseño, desarrollo, validación y eventual
mantenimiento del software.

Los destinatarios de este documento son el desarrollador del sistema,
el tutor académico, la gerencia del Bodegón Pa' Donde Natty C.A. y
los evaluadores del Trabajo Especial de Grado.

\subsection*{1.2 Alcance}

El sistema, denominado en adelante \textbf{``Sistema de Inventario IA''}
o simplemente \textbf{``el Sistema''}, es una aplicación web construida
sobre una arquitectura modular de tres capas (presentación, lógica de
negocio y persistencia) que permite digitalizar y automatizar el
control de inventarios del Bodegón Pa' Donde Natty C.A., incorporando
reconocimiento automatizado de productos mediante un modelo de
lenguaje multimodal de Inteligencia Artificial consumido a través de
su API.

El Sistema comprende autenticación y control de acceso, panel general
con indicadores clave, captura móvil vinculada por código QR,
reconocimiento de productos mediante IA, administración del catálogo,
gestión de stock y movimientos, alertas automáticas, reportes
gerenciales, respaldo y configuración. Quedan fuera del alcance la
facturación electrónica, contabilidad, comercio electrónico,
integración tributaria y gestión de recursos humanos.

\subsection*{1.3 Definiciones, acrónimos y abreviaturas}

\setstretch{1.0}
\noindent
\begin{tabular}{|p{3cm}|p{10.5cm}|}
\hline
\textbf{Término} & \textbf{Definición} \\
\hline
API & \textit{Application Programming Interface}; interfaz que permite la comunicación entre sistemas. \\
\hline
CRUD & \textit{Create, Read, Update, Delete}; operaciones básicas sobre registros. \\
\hline
ERD & \textit{Entity-Relationship Diagram}; modelo entidad-relación. \\
\hline
IA & Inteligencia Artificial. \\
\hline
IEEE & \textit{Institute of Electrical and Electronics Engineers}. \\
\hline
JWT & \textit{JSON Web Token}; mecanismo de autenticación basado en tokens firmados. \\
\hline
KPI & \textit{Key Performance Indicator}; indicador clave de desempeño. \\
\hline
MTR & Matriz de Trazabilidad de Requisitos. \\
\hline
PYME & Pequeña y Mediana Empresa. \\
\hline
QR & \textit{Quick Response}; código bidimensional de respuesta rápida. \\
\hline
REST & \textit{Representational State Transfer}; estilo arquitectónico para servicios web. \\
\hline
RF & Requerimiento Funcional. \\
\hline
RT & Requerimiento Técnico o no funcional. \\
\hline
SGBD & Sistema Gestor de Base de Datos. \\
\hline
SRS & \textit{Software Requirements Specification}; especificación de requerimientos de software. \\
\hline
3FN & Tercera Forma Normal. \\
\hline
UI / UX & \textit{User Interface / User Experience}; interfaz y experiencia de usuario. \\
\hline
\end{tabular}
\setstretch{1.5}

\subsection*{1.4 Referencias}

\setstretch{1.0}
\setlength{\parindent}{0pt}
\setlength{\parskip}{6pt}
IEEE Computer Society. (1998). \textit{IEEE Std 830-1998: IEEE
Recommended Practice for Software Requirements Specifications}.
Institute of Electrical and Electronics Engineers.

Sommerville, I. (2011). \textit{Ingeniería de software} (9a ed.).
Pearson Educación.

Pressman, R. y Maxim, B. (2020). \textit{Ingeniería del software: Un
enfoque práctico} (9a ed.). McGraw-Hill.

Date, C. J. (2001). \textit{Introducción a los sistemas de bases de
datos} (7a ed.). Pearson Educación.

Riera, G. (2026). \textit{Trabajo Especial de Grado: Sistema de
digitalización y registros de productos con uso de la Inteligencia
Artificial para control de inventarios en el Bodegón Pa' Donde Natty,
Ciudad Bolívar, Estado Bolívar}. Universidad Nororiental Privada Gran
Mariscal de Ayacucho.
\setstretch{1.5}
\setlength{\parindent}{1.25cm}
\setlength{\parskip}{0pt}

\section*{2. Descripción general}
\addcontentsline{toc}{section}{Anexo E. 2. Descripción general}

\subsection*{2.1 Perspectiva del producto}

El Sistema de Inventario IA es un producto de software nuevo,
construido desde cero, que reemplaza un proceso manual de control de
inventarios basado en conteos físicos, anotaciones en papel y hojas
de cálculo no integradas. La arquitectura se organiza en tres capas:
presentación con Next.js desplegada en Vercel; lógica de negocio con
Python y FastAPI desplegada en Render; y persistencia con PostgreSQL
gestionado mediante Supabase, que además provee sincronización en
tiempo real entre la PC y el teléfono.

\subsection*{2.2 Funciones del producto}

A alto nivel, el Sistema permite autenticar usuarios y restringir el
acceso por roles; vincular un teléfono mediante código QR; reconocer
productos automáticamente a partir de su imagen; registrar entradas y
salidas de mercancía; mantener catálogo de productos y categorías;
generar alertas de stock mínimo y vencimiento; consultar productos;
producir reportes gerenciales; presentar dashboard y respaldar la
información.

\subsection*{2.3 Características de los usuarios}

\setstretch{1.0}
\noindent
\begin{tabular}{|p{2.8cm}|p{4.5cm}|p{6.0cm}|}
\hline
\textbf{Perfil} & \textbf{Nivel técnico} & \textbf{Responsabilidades en el Sistema} \\
\hline
Administrador & Conocimientos básicos de informática & Gestión integral del sistema: usuarios, catálogo, configuración, respaldos y reportes gerenciales. \\
\hline
Gerente & Conocimientos básicos de informática & Consulta del panel general, lectura de reportes y supervisión de alertas. \\
\hline
Operador de almacén & Sin formación técnica avanzada & Captura móvil de productos, registro de entradas y salidas, confirmación del reconocimiento por IA. \\
\hline
\end{tabular}
\setstretch{1.5}

\subsection*{2.4 Restricciones}

El desarrollo y operación del Sistema están sujetos a restricciones
tecnológicas (Python/FastAPI, Next.js y PostgreSQL/Supabase), de
despliegue (Vercel, Render y Supabase en capa gratuita), de
conectividad, de terceros, legales y de tiempo.

\subsection*{2.5 Suposiciones y dependencias}

Se asume que la empresa cuenta con Internet estable, teléfono con
cámara y computadora con navegador moderno; que la API del modelo
multimodal mantiene condiciones operativas compatibles; y que la
gerencia facilita la información y participación del personal.

\section*{3. Requerimientos específicos}
\addcontentsline{toc}{section}{Anexo E. 3. Requerimientos específicos}

\subsection*{3.1 Requerimientos funcionales}

\setstretch{1.0}
\small
\begin{longtable}{|p{2.1cm}|p{4.1cm}|p{7.1cm}|}
\hline
\textbf{Código} & \textbf{Nombre} & \textbf{Especificación resumida} \\
\hline
RF-01 & Reconocimiento y registro de productos por imagen & El Sistema reconocerá productos desde imágenes JPEG/PNG capturadas con teléfono vinculado por QR, enviará la imagen a la API multimodal, validará la respuesta, asociará el producto con el catálogo, registrará el movimiento y actualizará el stock. Prioridad: Alta. \\
\hline
RF-02 & Registro de entradas y salidas & Permitirá registrar movimientos de entrada y salida con producto, cantidad, fecha, usuario y observaciones; recalculará el stock y sincronizará la información en tiempo real. Prioridad: Alta. \\
\hline
RF-03 & Alertas de stock mínimo y vencimiento & Generará alertas cuando el stock alcance el umbral mínimo o cuando un producto esté próximo a vencer, mostrando notificaciones en el dashboard. Prioridad: Alta. \\
\hline
RF-04 & Consulta y búsqueda de productos & Permitirá buscar por nombre, categoría o código, con filtros por estado y existencias, mostrando listados paginados sin modificar la base de datos. Prioridad: Alta. \\
\hline
RF-05 & Reportes gerenciales & Generará reportes de existencias, rotación y movimientos con KPIs, filtros por fecha, categoría, producto o usuario, y opción de exportación. Prioridad: Media. \\
\hline
RF-06 & Administración del catálogo & Permitirá crear, consultar, modificar y eliminar productos y categorías, validando campos obligatorios y unicidad de registros. Prioridad: Media. \\
\hline
RF-07 & Autenticación y control de acceso & Autenticará usuarios mediante credenciales, verificará contraseñas almacenadas con \textit{hash}, emitirá token JWT y restringirá acceso por rol. Prioridad: Alta. \\
\hline
RF-08 & Dashboard con indicadores y alertas & Presentará tarjetas resumen, alertas y gráficas con total de productos, stock crítico, productos por vencer y movimientos del día. Prioridad: Media. \\
\hline
RF-09 & Respaldo y configuración & Permitirá ejecutar respaldos descargables y configurar parámetros operativos como umbrales, días previos al vencimiento y sensibilidad de reconocimiento. Prioridad: Media. \\
\hline
\end{longtable}
\normalsize
\setstretch{1.5}

\subsection*{3.2 Requerimientos no funcionales o técnicos}

\begin{itemize}[leftmargin=*]
  \item \textbf{RT-01: Plataforma y backend.} El backend se desarrollará en Python con FastAPI, integrando el modelo multimodal mediante API y exponiendo servicios REST documentados con OpenAPI/Swagger. Prioridad: Alta.
  \item \textbf{RT-02: Base de datos relacional.} La persistencia empleará PostgreSQL gestionado por Supabase, con modelo normalizado hasta Tercera Forma Normal (3FN). Prioridad: Alta.
  \item \textbf{RT-03: Tiempo de respuesta.} El reconocimiento deberá responder en promedio en tres (03) segundos o menos, incluyendo latencia de red. Prioridad: Alta.
  \item \textbf{RT-04: Precisión.} La exactitud del reconocimiento deberá ser igual o superior al 90\,\%, con mecanismo de corrección manual. Prioridad: Alta.
  \item \textbf{RT-05: Interfaz y operabilidad.} La aplicación será operable desde navegador de PC y vista móvil por QR, con interfaz sencilla para personal sin formación técnica avanzada. Prioridad: Alta.
  \item \textbf{RT-06: Respaldo y configuración.} El Sistema dispondrá de respaldo bajo demanda y configuración de parámetros operativos. Prioridad: Media.
  \item \textbf{RT-07: Despliegue y tiempo real.} El frontend se desplegará en Vercel, el backend en Render y la base de datos en Supabase, procurando sincronización inmediata sin recarga manual. Prioridad: Alta.
\end{itemize}

\subsection*{3.3 Interfaces externas}

El Sistema presentará inicio de sesión, dashboard, vinculación por QR,
vista móvil de captura, resultado del reconocimiento, catálogo,
movimientos, consulta de inventario, reportes gerenciales y
configuración/respaldo. Requiere computadora con navegador moderno,
teléfono Android o iOS con cámara, API multimodal de IA sobre
HTTPS/JSON, API REST interna de FastAPI y servicios de Supabase.

\subsection*{3.4 Requerimientos de datos}

\setstretch{1.0}
\noindent
\begin{tabular}{|p{3cm}|p{10.5cm}|}
\hline
\textbf{Entidad} & \textbf{Atributos principales} \\
\hline
Usuario & id, nombre, correo, contraseña (hash), rol, estado, fecha\_creacion. \\
\hline
Categoria & id, nombre, descripcion, fecha\_creacion. \\
\hline
Producto & id, nombre, id\_categoria (FK), unidad\_medida, stock\_actual, stock\_minimo, precio, fecha\_vencimiento, imagen\_url, estado. \\
\hline
Movimiento & id, id\_producto (FK), id\_usuario (FK), tipo (entrada/salida), cantidad, fecha\_hora, observaciones. \\
\hline
Alerta & id, id\_producto (FK), tipo (stock\_minimo / vencimiento), fecha\_generacion, estado (activa/atendida). \\
\hline
\end{tabular}
\setstretch{1.5}

\section*{4. Matriz de trazabilidad de requisitos}
\addcontentsline{toc}{section}{Anexo E. 4. Matriz de trazabilidad de requisitos}

La Matriz de Trazabilidad de Requisitos (MTR) vincula cada
requerimiento funcional con su componente arquitectónico, las entidades
de base de datos involucradas y la pantalla de interfaz que lo
materializa.

\setstretch{1.0}
\scriptsize
\noindent
\begin{tabular}{|c|p{3.4cm}|p{3.0cm}|p{2.6cm}|c|}
\hline
\textbf{RF} & \textbf{Componente arquitectónico} & \textbf{Entidades BD} & \textbf{Pantalla UI} & \textbf{Trazado} \\
\hline
RF-01 & Reconocimiento IA + API & Producto, Movimiento & Captura móvil, Resultado & Sí \\
\hline
RF-02 & Gestión de stock (FastAPI) & Producto, Movimiento & Movimientos & Sí \\
\hline
RF-03 & Motor de alertas (FastAPI) & Producto, Alerta & Dashboard & Sí \\
\hline
RF-04 & Consulta (FastAPI) & Producto, Categoría & Inventario & Sí \\
\hline
RF-05 & Reportes (FastAPI/Next.js) & Movimiento, Producto & Reportes & Sí \\
\hline
RF-06 & Catálogo (FastAPI) & Producto, Categoría & Catálogo & Sí \\
\hline
RF-07 & Autenticación (FastAPI) & Usuario & Inicio de sesión & Sí \\
\hline
RF-08 & Agregación (FastAPI/Next.js) & Producto, Movimiento, Alerta & Dashboard & Sí \\
\hline
RF-09 & Respaldo/configuración (FastAPI) & Usuario, Producto & Configuración & Sí \\
\hline
\end{tabular}
\normalsize
\setstretch{1.5}

\vspace{8pt}

La cobertura alcanzada es del \textbf{100\,\%}: los nueve (09)
requerimientos funcionales poseen un componente arquitectónico, una o
más entidades de base de datos y una pantalla de interfaz asignados,
sin requerimientos huérfanos ni componentes injustificados.

\section*{Aprobación del documento}
\addcontentsline{toc}{section}{Anexo E. Aprobación del documento}

La presente Especificación de Requerimientos de Software (SRS) ha sido
revisada y aprobada como base formal para el desarrollo del Sistema de
Inventario IA del Bodegón Pa' Donde Natty C.A.

\vspace{1.2cm}

\noindent
\begin{minipage}{0.45\textwidth}
\centering
\rule{5.5cm}{0.4pt}\\[4pt]
\textbf{Autor}\\
Gabriel Riera\\
C.I.: V--31.153.478
\end{minipage}\hfill
\begin{minipage}{0.45\textwidth}
\centering
\rule{5.5cm}{0.4pt}\\[4pt]
\textbf{Tutor Académico}\\
Ing. Carlos Moreno
\end{minipage}

\vspace{1cm}

\begin{center}
Ciudad Bolívar, Mayo 2026.
\end{center}

% ============================================================


\end{document}